\documentclass[a4paper,11pt]{article}

% This file stands as a figure-only dummy of the paper, and serves to generate
% its more complicated figures with tikzexternal so that the results can then
% simply be \includegraphics'd into the paper.

% It relies on plot data generated with
%  - The main numeric implementation with command 'plot images' (image_Imt*.dat)
%  - Ditto with command 'plot E' through symlog.py (E*.dat, E*_symlog.tex)
%  - Ditto with command 'plot Pi' (Pi*.dat)
%  - Ditto with command 'plot J' (Jbub*.dat)
%  - regions.gp through join_contours.py (im*.dat)
%  - generate_plots.sage through trim_plots.sh (all the png's)


\usepackage{jheppub}

\usepackage[utf8]{inputenc}
\usepackage[american]{babel}

\usepackage{cleveref}
\usepackage{mathtools} % coloneq, allowdisplaybreaks

\usepackage{subcaption}
\usepackage{pgf,tikz}
\usetikzlibrary{%
    calc,
    positioning,
    patterns,
    decorations.pathreplacing,
    decorations.markings,
    decorations.pathmorphing,
    decorations.shapes,
    shapes.geometric,
    arrows.meta,
    external%
}

\usepackage{pgfplots}
\pgfplotsset{compat=1.18}
\usepgfplotslibrary{fillbetween}

% Settings for drawing the diagrams in TikZ

\newcommand{\tadpoleangle}{130}
\newlength{\tadpolesize}\setlength{\tadpolesize}{1cm}
\newlength{\bubblesize}\setlength{\bubblesize}{.33cm}
\newcommand{\bubbleaspect}{.8}
\newlength{\photonlength}\setlength{\photonlength}{.8cm}
\newcommand{\diagramscale}{1.25}
\newlength{\vertsize}\setlength{\vertsize}{2.25pt}
\tikzset{
    Ultra thick/.style={line width=3pt},
    miter cap/.style={{Triangle Cap[]}-{Triangle Cap[]}},
    pion/.style = { very thick },
    photon/.style = { thick, decorate, decoration={snake,segment length=.26\photonlength, amplitude=+.05\photonlength} },
    notoph/.style = { thick, decorate, decoration={snake,segment length=.26\photonlength, amplitude=-.05\photonlength} },
    diagrcut/.style={red,decorate,decoration={zigzag,segment length=1mm,amplitude=.25mm}},
    %
    bubble2/.style = {xscale=.8, transform shape},
    elbbub2/.style = {xscale=1.25, transform shape},
    bubble3/.style = {xscale=.64, transform shape},
    threeloop/.style = {scale=1.2, transform shape},
    %
    NLO/.pic = {
        \fill (0,0) circle[radius=\vertsize];
    },
    NNLO/.pic = {
        \fill (-\vertsize,-\vertsize) rectangle (\vertsize, \vertsize);
    },
    NNNLO/.pic = {
        \fill (90:2\vertsize) -- (210:2\vertsize) -- (330:2\vertsize) -- cycle;
    },
    %
    pics/tadpole/.style = {
        code = {
            \draw[pion] (0,0) .. controls (\tadpoleangle:\tadpolesize) and (180-\tadpoleangle:\tadpolesize) .. (0,0)
                foreach \i in {#1} { node[pos=.\i] (node\i) {} }
                ;
        }
    },
    pics/bubble/.style = {
        code = {
            \draw[pion] (0,0) ellipse [x radius=\bubblesize, y radius=\bubbleaspect\bubblesize];
            \foreach \i in {#1} {
                \coordinate (node\i) at
                (canvas polar cs:angle=\i, x radius=\bubblesize, y radius=\bubbleaspect\bubblesize) {};
            }
        }
    },
    pics/bubble/.default = {0}, % For some reason tadpole is fine with empty, but bubble needs default
    bubblejoint/.pic = {
        \fill[white] (-\bubblesize,-\bubblesize) rectangle (\bubblesize,\bubblesize);
        \draw[pion] (-\bubblesize,+\bubbleaspect\bubblesize) .. controls (0,+\bubbleaspect\bubblesize) and (0,-\bubbleaspect\bubblesize) .. (+\bubblesize,-\bubbleaspect\bubblesize);
        \draw[pion] (+\bubblesize,+\bubbleaspect\bubblesize) .. controls (0,+\bubbleaspect\bubblesize) and (0,-\bubbleaspect\bubblesize) .. (-\bubblesize,-\bubbleaspect\bubblesize);
    },
}

\newcounter{diagram}
\newenvironment{diagram}[1][]{%
%         \stepcounter{diagram}\arabic{diagram}
        \begin{tikzpicture}[scale=\diagramscale, transform shape, baseline={(0,0)},#1]%
        \path[use as bounding box] (-\photonlength,-\bubblesize) rectangle (+\photonlength,+\bubblesize);
    }{%
        % Smash vertically, keep horizontal extent
        \end{tikzpicture}%
    }

    
% P. Tol's colorblind-friendly plot colors
\definecolor{plotI}{HTML}{0077BB}
\definecolor{plotII}{HTML}{33BBEE}
\definecolor{plotIII}{HTML}{009988}
\definecolor{plotIV}{HTML}{EE7733}
\definecolor{plotV}{HTML}{CC3311}
\definecolor{plotVI}{HTML}{DDAA33}
\definecolor{plotVII}{HTML}{AA3377}

% Colors for various kinematic regions
\colorlet{spacelike}{plotI}
\colorlet{subthr}{plotIV}
\colorlet{twopi}{plotIII}
\colorlet{fourpi}{plotVII}
\colorlet{shaded}{gray!30}

% Additional tikz settings new to the numerics paper

\usepackage{etoolbox}

% These are from the bounds paper - colorblind friendly!
\definecolor{plotI}{HTML}{0077BB}
\definecolor{plotII}{HTML}{33BBEE}
\definecolor{plotIII}{HTML}{009988}
\definecolor{plotIV}{HTML}{EE7733}
\definecolor{plotV}{HTML}{CC3311}
\definecolor{plotVI}{HTML}{DDAA33}
\definecolor{plotVII}{HTML}{AA3377}

% These are from the Kdf papers - look good on a 22pt legend
\tikzset{
    dash0/.style={solid},
    dash1/.style={dash pattern=on 8pt off 2pt on 2pt off 2pt},
    dash2/.style={dash pattern=on 6pt off 2pt on 2pt off 2pt on 2pt off 2pt},
    dashA/.style={dash pattern=on 4.4pt off 4.4pt},
    dashB/.style={dash pattern=on 2pt off 2pt},
    dashC/.style={dash pattern=on 1.2pt off 1pt},
    dash3/.style={dash pattern=on 5.5pt off 1.5pt on 1.5pt off 1.5pt on 1.5pt off 1.5pt on 1.5pt off 1.5pt},
    dash4/.style={dash pattern=on 5.2pt off 1.2pt on 1.2pt off 1.2pt on 1.2pt off 1.2pt on 1.2pt off 1.2pt on 1.2pt off 1.2pt},
    dash5/.style={dash pattern=on 5.2pt off 0.9pt on 0.9pt off 0.9pt on 0.9pt off 0.9pt on 0.9pt off 0.9pt on 0.9pt off 0.9pt},
    %
    line0/.style={plotI,   dash0},
    line1/.style={plotII,  dash1},
    line2/.style={plotIII, dash2},
    lineA/.style={plotIV,  dashA},
    lineB/.style={plotV,   dashB},
    line3/.style={plotVI,  dash3},
    line4/.style={plotVII, dash4},
    line5/.style={plotVII, dash5},
    %
    real/.style={line0, thick},
    imag/.style={lineA, thick},
    %
    bessel/.style={dashA, plotVII, thick},
    polylog/.style={dash2, plotVI, thick},
    eisen/.style={dash0, plotIII, thick},
    poisson/.style={dash1, plotIV, thick},
    series dash/.style={dashC, thick},
    series/.style={series dash, plotV},
    secdec/.style={dashB, plotI, thick, join=round},
    feyntrop/.style={dash4, plotII, thick},
    amflow/.style={solid, thick, mark=+, draw=none, /pgfplots/error bars/y dir=both, /pgfplots/error bars/y explicit},
    }

\colorlet{E1}{plotI}
\colorlet{E2}{plotIII}
\colorlet{E3}{plotV}
\colorlet{E4}{plotIV}
\colorlet{E5}{plotVI}
\colorlet{E6}{plotVII}
\colorlet{realextra}{plotIV}

\tikzset{
    branchcut/.style={decorate, decoration={zigzag, segment length=1mm, amplitude=.4mm}, rounded corners=.01pt},
    realpos/.style={thick, plotV},
    realneg/.style={thick, plotII},
    nocut/.style={densely dotted},
    C0/.style={-{>[scale=.7]}},
    C1/.style={{<[scale=.7]}-{>[scale=.7]}},
    realline/.style={line width=2pt, miter cap},
    realline shadow/.style={white, dashed, thick},
    }

% Sub- and above-threshold arc radii
\def\subrad{0.16666667} % = 1/6
\def\suprad{0.28867513} % = sqrt(3)/6

% Arguments: (options) center radius
\newcommand{\addqarc}[3][]{%
    \addplot[#1] (%
        {exp(-2*pi*(#3)*sin(x))*cos(360*(#2+(#3)*cos(x)))},%
        {exp(2*pi*-(#3)*sin(x))*sin(360*(#2+(#3)*cos(x)))});}
\newcommand{\addqarcs}[3][]{\addqarc[#1]{#2}{#3}\addqarc[#1]{-(#2)}{#3}}

\newcommand{\qplot}[1]{
    \begin{subfigure}[t]{.3\columnwidth}
        \begin{tikzpicture}
            \begin{axis}[
                width=\columnwidth, height=\columnwidth, scale only axis, axis on top,
                xmin=-1, xmax=1,
                ymin=-1, ymax=1,
                ticks=none, axis lines=none
                ]
                \addplot graphics [
                    xmin=-1, xmax=1,
                    ymin=-1, ymax=1
                ] {\plotdir/#1};
                \addqarc[realline shadow, domain=60:180]{\subrad}{\subrad}
                \addqarc[realline shadow, domain=0:120]{-\subrad}{\subrad}
                \addqarc[realline shadow, domain=90:150]{.5}{\suprad}
                \addqarc[realline shadow, domain=30:90]{-.5}{\suprad}
            \end{axis}
        \end{tikzpicture}
    \end{subfigure}}

\pgfplotsset{
    every axis/.style={set layers},
    tcontour/.style={very thin, gray, mark=none, domain=-100:100, restrict y to domain=#1}}

% The remainder is for the result plots.
% Preliminary versions with more variety can be found in test_plots_v2.tex

\tikzset{
    threshold/.style={dash pattern=on 0pt off 2pt, cap=round, line width=1.5pt, opacity=.6},
%     imdash/.style={dash pattern=on 3pt off 2.5pt},
%     series dotted/.style={dash pattern=on 1pt off 1pt},
    real/.style={plotI, very thick, dash0},
    imag/.style={plotIV, very thick, dashA},
%     series/.style={real, thick, series dotted},
    ours/.style={real, very thick, solid},
%     secdec/.style={real, thick, dash pattern=on 3pt off 1pt},
%     feyntrop/.style={real, thick, dash pattern=on 1pt off 2pt},
    errorband/.style={solid, draw=none, opacity=.35},
    faint errorband/.style={opacity=.35},
    dark errorband/.style={opacity=.6},
    total/.style={line width=2pt, black}}

\newlength{\legendwidth}\setlength{\legendwidth}{22pt}
\pgfplotsset{
    complex legend/.style={legend image code/.code={
        \draw [real, #1] (0pt,.5ex) -- (\legendwidth,.5ex);
        \draw [imag, #1] (0pt,-.2ex) -- (\legendwidth,-.2ex);
        }},
%     complex legend with PSD/.style={legend image code/.code={
%         \fill [#1, draw=none, secdec] (0pt,-.5ex) rectangle (\legendwidth, 1ex);
%         \draw [#1, real] (0pt,.5ex) -- (\legendwidth,.5ex);
%         \draw [#1, imag] (0pt,0ex) -- (\legendwidth,0ex);
%         }},
    }
\pgfplotsset{
    legend with error/.style={legend image code/.code={
        \fill [#1, solid, errorband] (0pt,+.7ex) rectangle (\legendwidth,-.3ex);
        \draw [#1]            (0pt,+.2ex) -- (\legendwidth,+.2ex);
        }},
    legend with different error/.style 2 args={legend image code/.code={
        \fill [#2, errorband] (0pt,+.7ex) rectangle (\legendwidth,-.3ex);
        \draw [#1]            (0pt,+.2ex) -- (\legendwidth,+.2ex);
        }},
    double legend/.style={legend image code/.code={
        \draw [ours,     #1] (0pt,+.7ex) -- (\legendwidth,+.7ex);
        \draw [secdec,   #1] (0pt,+.1ex) -- (\legendwidth,+.1ex);
        }},
    triple legend/.style={legend image code/.code={
        \draw [ours,     #1] (0pt,+.7ex) -- (\legendwidth,+.7ex);
        \draw [secdec,   #1] (0pt,+.1ex) -- (\legendwidth,+.1ex);
        \draw [feyntrop, #1] (0pt,-.5ex) -- (\legendwidth,-.5ex);
        }},
    }

\newcommand{\plotE}[5]{%
    \begin{tikzpicture}
        \def\Eplotfile{\plotdir/prec64/E#1}
        \def\Eplotlabel{#2}
        \def\Eplotymin{#3}
        \def\Eplotymax{#4}
        \def\Eplotthreshold{#5}
        \begin{axis}[
                width=\Eplotwidth, height=\Eplotheight, scale only axis,
                ylabel={$\Eplotlabel$},
                xmin=\Eplotxmin, xmax=\Eplotxmax,
                xtick={0,4,16}, xmajorgrids=true,
                minor xtick={-9,-8,...,23},
                xticklabels={},
                ymin=\Eplotymin, ymax=\Eplotymax,
                unbounded coords=jump
                ]
            \drawzerolines{\Eplotxmin}{\Eplotxmax}{\Eplotymin}{\Eplotymax}

            \addplot[real]   table[x=t, y=Re, col sep=tab, restrict y to domain=\Eplotymin:\Eplotymax] {\Eplotfile e.dat};
            \addplot[imag]   table[x=t, y=Im, col sep=tab, restrict x to domain=\Eplotthreshold:\Eplotxmax, restrict y to domain=\Eplotymin:\Eplotymax] {\Eplotfile e.dat};

%             \addplot[series] table[x=t, y=Re, col sep=tab, restrict y to domain=\Eplotymin:\Eplotymax] {\Eplotfile 0.dat};
            \addplot[real, series dash] table[x=t, y=Re, col sep=tab, restrict y to domain=\Eplotymin:\Eplotymax] {\Eplotfile 08.dat};
%             \addplot[series] table[x=t, y=Re, col sep=tab, restrict y to domain=\Eplotymin:\Eplotymax] {\Eplotfile 016.dat};

            \addplot[amflow, real] table[x=t, y=Re, y error plus expr=abs(\thisrow{MaxRe}-\thisrow{Re}), y error minus expr=abs(\thisrow{MinRe}-\thisrow{Re})] {\Eplotfile a.dat};
            \addplot[amflow, imag, restrict x to domain=\Eplotthreshold:\Eplotxmax] table[x=t, y=Im, y error plus expr=abs(\thisrow{MaxIm}-\thisrow{Im}), y error minus expr=abs(\thisrow{MinIm}-\thisrow{Im})] {\Eplotfile a.dat};

        \end{axis}

        \begin{semilogyaxis}[
                shift={(0,-(\doublenormEplotheight + \Eplotsep))},
                width=\Eplotwidth, height=\doublenormEplotheight, scale only axis,
                ylabel={$\delta\Eplotlabel$},
                xmin=\Eplotxmin, xmax=\Eplotxmax,
                xtick={0,4,16}, xmajorgrids=true,
                minor xtick={-9,-8,...,23},
                xlabel={\xlabel},
                ymin=\Eplotlogymin, ymax=\Eplotlogymax,
                ytick={1e-3,1e-6,1e-9,1e-12,1e-15},
                minor ytick={1e-1,1e-2,1e-4,1e-5,1e-7,1e-8,1e-10,1e-11,1e-13,1e-14,1e-16,1e-17},
                ymajorgrids=true,
                unbounded coords=jump
                ]

            \addplot[draw=none, name path=zero, domain=\Eplotxmin:\Eplotxmax, samples=2] {\Eplotlogymin};
            \addplot[draw=none, name path=psd] table[x=t, y=AbsRe, col sep=tab] {\Eplotfile s.dat};
            \ifnumless{#1}{5}{%
                \addplot[gray, errorband] fill between[of=zero and psd];

                \addplot[draw=none, name path=ft] table[x=t, y=AbsRe, col sep=tab] {\Eplotfile f.dat};
                \addplot[gray, faint errorband] fill between[of=zero and ft];%
                }{%
                \addplot[gray, dark errorband] fill between[of=zero and psd];
                }
            \addplot[real]   table[x=t, y=AbsRe, col sep=tab] {\Eplotfile e.dat};
            \addplot[imag]   table[x=t, y=AbsIm, col sep=tab] {\Eplotfile e.dat};
%             \addplot[series] table[x=t, y=AbsRe, col sep=tab] {\Eplotfile 0.dat};
            \addplot[real, series dash] table[x=t, y=AbsRe, col sep=tab] {\Eplotfile 08.dat};
%             \addplot[series] table[x=t, y=AbsRe, col sep=tab] {\Eplotfile 016.dat};


        \end{semilogyaxis}
        \useasboundingbox (0,-\doublenormEplotheight) + (0,-\Eplotsep) + (0,-.5cm);  % Force larger bounding box
    \end{tikzpicture}}


\newcommand{\plotEsymlog}[3]{%
    \begin{tikzpicture}
        \def\Eplotfile{\plotdir/prec64/E#1}
        \def\Eplotlabel{#2}
        \def\Eplotthreshold{#3}
        \def\symlog{_symlog-15-1}
        \input{\plotdir/prec64/E#1e\symlog.tex}
        \begin{axis}[
            width=\Eplotwidth, height=\symlogEplotheight, scale only axis,
            xmin=\Eplotxmin, xmax=\Eplotxmax, xlabel={},
            xtick={0,4,16}, xmajorgrids=true, xticklabels={},
            minor xtick={-9,-8,...,23},
            ymin=\symlogymin, ymax=\symlogymax, ylabel={\small$\Re\Eplotlabel$},
            ytick=\symlogsparseytick,
            % This is \symlogsparseyticklabels, further pruned
            yticklabels={,$-10^{-6}$,,$-10^{-12}$,,$0$,,$+10^{-12}$,,$+10^{-6}$,},
            yticklabel style={font=\tiny},
            ymajorgrids=true, yminorgrids=true
            ]
            \addplot[gray, amflow] table[x=t, y=NormRe, y error plus expr=abs(\thisrow{NormMaxRe}-\thisrow{NormRe}), y error minus expr=abs(\thisrow{NormMinRe}-\thisrow{NormRe})] {\Eplotfile a\symlog.dat};

            \addplot[eisen] table[x=t, y=NormRe] {\Eplotfile e\symlog.dat};
            \addplot[draw=none, name path=upper] table[x=t, y=NormMaxRe] {\Eplotfile e\symlog.dat};
            \addplot[draw=none, name path=lower] table[x=t, y=NormMinRe] {\Eplotfile e\symlog.dat};
            \addplot[eisen, errorband] fill between[of=lower and upper];

            \addplot[secdec] table[x=t, y=NormRe] {\Eplotfile s\symlog.dat};
            \addplot[draw=none, name path=upper] table[x=t, y=NormMaxRe] {\Eplotfile s\symlog.dat};
            \addplot[draw=none, name path=lower] table[x=t, y=NormMinRe] {\Eplotfile s\symlog.dat};
            \addplot[secdec, errorband] fill between[of=lower and upper];
            \ifnumless{#1}{5}{%
                \begin{scope}[/pgfplots/restrict x to domain=\Eplotxmin:\Eplotthreshold]
                    \addplot[feyntrop] table[x=t, y=NormRe] {\Eplotfile f\symlog.dat};
                    \addplot[draw=none, name path=upper] table[x=t, y=NormMaxRe] {\Eplotfile f\symlog.dat};
                    \addplot[draw=none, name path=lower] table[x=t, y=NormMinRe] {\Eplotfile f\symlog.dat};
                    \addplot[feyntrop, errorband] fill between[of=lower and upper];%
                \end{scope}
                }{}

            \addplot[series, restrict x to domain=-\Eplotthreshold:\Eplotthreshold] table[x=t, y=NormRe] {\Eplotfile 0\symlog.dat};
            \addplot[draw=none, name path=upper, restrict x to domain=-\Eplotthreshold:\Eplotthreshold] table[x=t, y=NormMaxRe] {\Eplotfile 0\symlog.dat};
            \addplot[draw=none, name path=lower, restrict x to domain=-\Eplotthreshold:\Eplotthreshold] table[x=t, y=NormMinRe] {\Eplotfile 0\symlog.dat};
            \addplot[series, errorband] fill between[of=lower and upper];

        \end{axis}

        \begin{axis}[
            shift={(0,-(\symlogEplotheight + \Eplotsep))},
            width=\Eplotwidth, height=\symlogEplotheight, scale only axis,
            xmin=\Eplotxmin, xmax=\Eplotxmax, xlabel={\xlabel},
            xtick={0,4,16}, xmajorgrids=true,
            minor xtick={-9,-8,...,23},
            ymin=\symlogymin, ymax=\symlogymax, ylabel={\small$\Im\Eplotlabel$},
            ytick=\symlogsparseytick,
            % This is \symlogsparseyticklabels, further pruned
            yticklabels={,$-10^{-6}$,,$-10^{-12}$,,$0$,,$+10^{-12}$,,$+10^{-6}$,},
            yticklabel style={font=\tiny},
            ymajorgrids=true, yminorgrids=true,
            axis background/.style={%
                preaction={
                    path picture={
                        \path[pattern=north east lines, pattern color=gray!50] (axis cs:\Eplotxmin,\symlogymax) rectangle (axis cs:\Eplotthreshold,\symlogymin);
                    }}}
            ]
            \addplot[gray, amflow, restrict x to domain=\Eplotxmin:\Eplotthreshold] table[x=t, y=Im, y error plus expr=abs(\thisrow{MaxIm}-\thisrow{Im}), y error minus expr=abs(\thisrow{MinIm}-\thisrow{Im})] {\Eplotfile a\symlog.dat};
            \addplot[gray, amflow, restrict x to domain=\Eplotthreshold:\Eplotxmax] table[x=t, y=NormIm, y error plus expr=abs(\thisrow{NormMaxIm}-\thisrow{NormIm}), y error minus expr=abs(\thisrow{NormMinIm}-\thisrow{NormIm})] {\Eplotfile a\symlog.dat};

            \addplot[eisen, restrict x to domain=\Eplotxmin:\Eplotthreshold] table[x=t, y=Im] {\Eplotfile e\symlog.dat};
            \addplot[draw=none, name path=upper, restrict x to domain=\Eplotxmin:\Eplotthreshold] table[x=t, y=MaxIm] {\Eplotfile e\symlog.dat};
            \addplot[draw=none, name path=lower, restrict x to domain=\Eplotxmin:\Eplotthreshold] table[x=t, y=MinIm] {\Eplotfile e\symlog.dat};
            \addplot[eisen, errorband] fill between[of=lower and upper];

            \addplot[secdec, restrict x to domain=\Eplotxmin:\Eplotthreshold] table[x=t, y=Im] {\Eplotfile s\symlog.dat};
            \addplot[draw=none, name path=upper, restrict x to domain=\Eplotxmin:\Eplotthreshold] table[x=t, y=MaxIm] {\Eplotfile s\symlog.dat};
            \addplot[draw=none, name path=lower, restrict x to domain=\Eplotxmin:\Eplotthreshold] table[x=t, y=MinIm] {\Eplotfile s\symlog.dat};
            \addplot[secdec, errorband] fill between[of=lower and upper];
            \ifnumless{#1}{5}{%
                \addplot[feyntrop, restrict x to domain=\Eplotxmin:\Eplotthreshold] table[x=t, y=NormRe] {\Eplotfile f\symlog.dat};
                \addplot[draw=none, name path=upper, restrict x to domain=\Eplotxmin:\Eplotthreshold] table[x=t, y=NormMaxRe] {\Eplotfile f\symlog.dat};
                \addplot[draw=none, name path=lower, restrict x to domain=\Eplotxmin:\Eplotthreshold] table[x=t, y=NormMinRe] {\Eplotfile f\symlog.dat};
                \addplot[feyntrop, faint errorband] fill between[of=lower and upper];%
                }{}

            \addplot[eisen, restrict x to domain=\Eplotthreshold:\Eplotxmax] table[x=t, y=NormIm] {\Eplotfile e\symlog.dat};
            \addplot[draw=none, name path=upper, restrict x to domain=\Eplotthreshold:\Eplotxmax] table[x=t, y=NormMaxIm] {\Eplotfile e\symlog.dat};
            \addplot[draw=none, name path=lower, restrict x to domain=\Eplotthreshold:\Eplotxmax] table[x=t, y=NormMinIm] {\Eplotfile e\symlog.dat};
            \addplot[eisen, errorband] fill between[of=lower and upper];

            \addplot[secdec, restrict x to domain=\Eplotthreshold:\Eplotxmax] table[x=t, y=NormIm] {\Eplotfile s\symlog.dat};
            \addplot[draw=none, name path=upper, restrict x to domain=\Eplotthreshold:\Eplotxmax] table[x=t, y=NormMaxIm] {\Eplotfile s\symlog.dat};
            \addplot[draw=none, name path=lower, restrict x to domain=\Eplotthreshold:\Eplotxmax] table[x=t, y=NormMinIm] {\Eplotfile s\symlog.dat};
            \addplot[secdec, dark errorband] fill between[of=lower and upper];


        \end{axis}
        \useasboundingbox (0,-\symlogEplotheight) + (0,-\Eplotsep) + (0,-.5cm); % Force larger bounding box
    \end{tikzpicture}
    }

\begin{comment}

\newcommand{\plotEsymlog}[3]{%
    \begin{tikzpicture}
        \def\Eplotfile{\plotdir/prec64/E#1}
        \def\Eplotlabel{#2}
        \def\Eplotthreshold{#3}
        \def\symlog{_symlog-15-1}
        \input{\plotdir/prec64/E#1e\symlog.tex}
        \begin{axis}[
            width=\Eplotwidth, height=\symlogEplotheight, scale only axis,
            xmin=\Eplotxmin, xmax=\Eplotxmax, xlabel={},
            xtick={0,4,16}, xmajorgrids=true, xticklabels={},
            minor xtick={-9,-8,...,23},
            ymin=\symlogymin, ymax=\symlogymax, ylabel={\small$\Re\Eplotlabel$},
            ytick=\symlogsparseytick,
            % This is \symlogsparseyticklabels, further pruned
            yticklabels={,$-10^{-6}$,,$-10^{-12}$,,$0$,,$+10^{-12}$,,$+10^{-6}$,},
            yticklabel style={font=\tiny},
            ymajorgrids=true, yminorgrids=true
            ]
            \addplot[real, series, restrict x to domain=-\Eplotthreshold:\Eplotthreshold] table[x=t, y=NormRe] {\Eplotfile 0\symlog.dat};

            \addplot[real, gray] table[x=t, y=NormRe] {\Eplotfile e\symlog.dat};
            \addplot[draw=none, name path=upper] table[x=t, y=NormMaxRe] {\Eplotfile e\symlog.dat};
            \addplot[draw=none, name path=lower] table[x=t, y=NormMinRe] {\Eplotfile e\symlog.dat};
            \addplot[gray, errorband] fill between[of=lower and upper];

            \addplot[real] table[x=t, y=NormRe] {\Eplotfile s\symlog.dat};
            \addplot[draw=none, name path=upper] table[x=t, y=NormMaxRe] {\Eplotfile s\symlog.dat};
            \addplot[draw=none, name path=lower] table[x=t, y=NormMinRe] {\Eplotfile s\symlog.dat};
            \ifnumless{#1}{5}{%
                \addplot[real, errorband] fill between[of=lower and upper];

                \begin{scope}[/pgfplots/restrict x to domain=\Eplotxmin:\Eplotthreshold]
                    \addplot[real, thin] table[x=t, y=NormRe] {\Eplotfile f\symlog.dat};
                    \addplot[draw=none, name path=upper] table[x=t, y=NormMaxRe] {\Eplotfile f\symlog.dat};
                    \addplot[draw=none, name path=lower] table[x=t, y=NormMinRe] {\Eplotfile f\symlog.dat};
                    \addplot[real, faint errorband] fill between[of=lower and upper];%
                \end{scope}
                }{%
                \addplot[real, dark errorband] fill between[of=lower and upper];
                }

            \addplot[real, amflow] table[x=t, y=NormRe, y error plus expr=abs(\thisrow{NormMaxRe}-\thisrow{NormRe}), y error minus expr=abs(\thisrow{NormMinRe}-\thisrow{NormRe})] {\Eplotfile a\symlog.dat};

        \end{axis}

        \begin{axis}[
            shift={(0,-(\symlogEplotheight + \Eplotsep))},
            width=\Eplotwidth, height=\symlogEplotheight, scale only axis,
            xmin=\Eplotxmin, xmax=\Eplotxmax, xlabel={\xlabel},
            xtick={0,4,16}, xmajorgrids=true,
            minor xtick={-9,-8,...,23},
            ymin=\symlogymin, ymax=\symlogymax, ylabel={\small$\Im\Eplotlabel$},
            ytick=\symlogsparseytick,
            % This is \symlogsparseyticklabels, further pruned
            yticklabels={,$-10^{-6}$,,$-10^{-12}$,,$0$,,$+10^{-12}$,,$+10^{-6}$,},
            yticklabel style={font=\tiny},
            ymajorgrids=true, yminorgrids=true,
            axis background/.style={%
                preaction={
                    path picture={
                        \path[pattern=north east lines, pattern color=gray!50] (axis cs:\Eplotxmin,\symlogymax) rectangle (axis cs:\Eplotthreshold,\symlogymin);
                    }}}
            ]

            \addplot[imag, gray, restrict x to domain=\Eplotxmin:\Eplotthreshold] table[x=t, y=Im] {\Eplotfile e\symlog.dat};
            \addplot[draw=none, name path=upper, restrict x to domain=\Eplotxmin:\Eplotthreshold] table[x=t, y=MaxIm] {\Eplotfile e\symlog.dat};
            \addplot[draw=none, name path=lower, restrict x to domain=\Eplotxmin:\Eplotthreshold] table[x=t, y=MinIm] {\Eplotfile e\symlog.dat};
            \addplot[imag, errorband, gray] fill between[of=lower and upper];

            \addplot[imag, restrict x to domain=\Eplotxmin:\Eplotthreshold] table[x=t, y=Im] {\Eplotfile s\symlog.dat};
            \addplot[draw=none, name path=upper, restrict x to domain=\Eplotxmin:\Eplotthreshold] table[x=t, y=MaxIm] {\Eplotfile s\symlog.dat};
            \addplot[draw=none, name path=lower, restrict x to domain=\Eplotxmin:\Eplotthreshold] table[x=t, y=MinIm] {\Eplotfile s\symlog.dat};
            \ifnumless{#1}{5}{%
                \addplot[imag, errorband] fill between[of=lower and upper];

                \addplot[imag, thin, restrict x to domain=\Eplotxmin:\Eplotthreshold] table[x=t, y=NormRe] {\Eplotfile f\symlog.dat};
                \addplot[draw=none, name path=upper, restrict x to domain=\Eplotxmin:\Eplotthreshold] table[x=t, y=NormMaxRe] {\Eplotfile f\symlog.dat};
                \addplot[draw=none, name path=lower, restrict x to domain=\Eplotxmin:\Eplotthreshold] table[x=t, y=NormMinRe] {\Eplotfile f\symlog.dat};
                \addplot[imag, faint errorband] fill between[of=lower and upper];%
                }{%
                \addplot[imag, dark errorband] fill between[of=lower and upper];
                }

            \addplot[draw=none, name path=upper, restrict x to domain=\Eplotthreshold:\Eplotxmax] table[x=t, y=NormMaxIm] {\Eplotfile e\symlog.dat};
            \addplot[draw=none, name path=lower, restrict x to domain=\Eplotthreshold:\Eplotxmax] table[x=t, y=NormMinIm] {\Eplotfile e\symlog.dat};
            \addplot[imag, errorband, gray] fill between[of=lower and upper];

            \addplot[imag, restrict x to domain=\Eplotthreshold:\Eplotxmax] table[x=t, y=NormIm] {\Eplotfile s\symlog.dat};
            \addplot[draw=none, name path=upper, restrict x to domain=\Eplotthreshold:\Eplotxmax] table[x=t, y=NormMaxIm] {\Eplotfile s\symlog.dat};
            \addplot[draw=none, name path=lower, restrict x to domain=\Eplotthreshold:\Eplotxmax] table[x=t, y=NormMinIm] {\Eplotfile s\symlog.dat};
            \addplot[imag, dark errorband] fill between[of=lower and upper];

            \addplot[imag, amflow, restrict x to domain=\Eplotxmin:\Eplotthreshold] table[x=t, y=Im, y error plus expr=abs(\thisrow{MaxIm}-\thisrow{Im}), y error minus expr=abs(\thisrow{MinIm}-\thisrow{Im})] {\Eplotfile a\symlog.dat};
            \addplot[imag, amflow, restrict x to domain=\Eplotthreshold:\Eplotxmax] table[x=t, y=NormIm, y error plus expr=abs(\thisrow{NormMaxIm}-\thisrow{NormIm}), y error minus expr=abs(\thisrow{NormMinIm}-\thisrow{NormIm})] {\Eplotfile a\symlog.dat};

        \end{axis}
        \useasboundingbox (0,-\symlogEplotheight) + (0,-\Eplotsep) + (0,-.5cm); % Force larger bounding box
    \end{tikzpicture}

\end{comment}

\newcommand{\plotTime}[1]{
    \def\plotfile{\plotdir/prec64/E#1}
    \ifnumless{#1}{5}{%
        \addplot[imag, ours  ,   E#1, triple legend={E#1}] table[x=t, y expr=\thisrow{Time}/1e9] {\plotfile e.dat};
        \addplot[imag, secdec,   E#1, forget plot        ] table[x=t, y expr=\thisrow{Time}/1e9] {\plotfile s.dat};
        \addplot[imag, feyntrop, E#1, forget plot        ] table[x=t, y expr=\thisrow{Time}/1e9] {\plotfile f.dat};
        }{%
        \addplot[imag, ours  ,   E#1, double legend={E#1}] table[x=t, y expr=\thisrow{Time}/1e9] {\plotfile e.dat};
        \addplot[imag, secdec,   E#1, forget plot        ] table[x=t, y expr=\thisrow{Time}/1e9] {\plotfile s.dat};
        }
        index = 0
    }


\newcommand{\plotJbub}[1]{%
    \addplot[real, E#1, complex legend={E#1}] table[x=t, y=ReJ#1, col sep=tab] {\plotdir/Jbub#1.dat};
    \addplot[imag, E#1, forget plot, restrict x to domain=4:\Eplotxmax] table[x=t, y=ImJ#1, col sep=tab] {\plotdir/Jbub#1.dat};
    }

\newcommand{\drawzerolines}[4]{
            \draw[thin, gray] (#1, 0) -- (#2, 0);
            \draw[thin, gray] (0, #3) -- (0, #4);
            }

\newcommand{\newsetlength}[2]{\newlength{#1}\setlength{#1}{#2}}
\newsetlength{\Eplotwidth}{6cm}
\newsetlength{\Eplotheight}{3cm}
\newsetlength{\Eplotsep}{1mm}
\newsetlength{\normEplotheight}{.8cm}
\newsetlength{\doublenormEplotheight}{2\normEplotheight}\addtolength{\doublenormEplotheight}{\Eplotsep}
\newsetlength{\symlogEplotheight}{.7\Eplotheight}

\newcommand{\Eplotxmin}{-8}
\newcommand{\Eplotxmax}{+24}
\newcommand{\Eplotlogymin}{1e-18}
\newcommand{\Eplotlogymax}{1}
\newcommand{\xlabel}{}

\newcommand{\samplelength}{23pt}
\newcommand{\sample}[1]{\tikz[baseline=-0.25ex]{
    \draw[real, secdec, #1] (0,0) -- (\samplelength,0);
    \draw[real, ours,   #1] (0,.5ex) -- (\samplelength,.5ex);
    }}
\newcommand{\Resample}{%
    blue, solid:%
    \tikzset{external/export next=false}
    \tikz[baseline=-0.5ex]{
        \draw[real] (0,0) -- (19.5pt,0);
        }}
\newcommand{\Imsample}{%
    orange, dashed:%
    \tikzset{external/export next=false}
    \tikz[baseline=-0.6ex]{
        \draw[imag] (0,0) -- (19.5pt,0);
        }}
\newcommand{\seriessample}{%
    blue, dotted:%
    \tikzset{external/export next=false}
    \tikz[baseline=-0.6ex]{
        \draw[real, series] (0,0) -- (\samplelength,0);
        }}

\newcommand{\linesample}[1]{%
    \tikzset{external/export next=false}
    \tikz[baseline=-0.6ex]{
        \draw[#1] (0,0) -- (19.5pt,0);
        }}

\newsetlength{\hybridplotwidth}{.8\Eplotwidth}
\newsetlength{\hybridplotheight}{\symlogEplotheight}

\tikzset{
    totval/.style={plotI, thick, solid},
    headval/.style={plotII, thick, densely dashed},
    tailval/.style={plotIII, thick, densely dotted},
    }


% From https://tex.stackexchange.com/questions/58548/is-it-possible-to-change-the-color-of-a-single-bar-when-the-bar-plot-is-based-on
\pgfplotsset{
    discard if/.style 2 args={
        x filter/.code={
            \edef\tempa{\thisrow{#1}}
            \edef\tempb{#2}
            \ifx\tempa\tempb
                \def\pgfmathresult{inf}
            \fi
        }
    },
    discard if not/.style 2 args={
        x filter/.code={
            \edef\tempa{\thisrow{#1}}
            \edef\tempb{#2}
            \ifx\tempa\tempb
            \else
                \def\pgfmathresult{inf}
            \fi
        }
    }
}

\newcommand{\plothybrid}[9]{%
    \def\hybridfile{\plotdir/hybrid_err_t-2.0.dat}
    \def\plainname{#1}
    \def\prettyname{#2}
    \def\chimin{#3}
    \def\chiref{#4}
    \def\chimax{#5}
    \def\relmin{#6}
    \def\relmax{#7}
    \def\valatchiref{#8}
    \def\errmin{#9}
    \begin{tikzpicture}
        \begin{semilogxaxis}[
            width=\hybridplotwidth, height=\hybridplotheight, scale only axis,
            xmin=\chimin, xmax=\chimax,
%             ymin=\relmin, ymax=\relmax,
            xticklabels={},
            ylabel={$\prettyname (\text{relative})$}
%             legend pos=north east, reverse legend
            ]
            \addplot[mark=*] coordinates { (\chiref, 1) };
%             \addplot[tailval] table[x=chi, y expr={\thisrow{tail}/\valatchiref}, discard if not={source}{\plainname}, col sep=tab] {\hybridfile};
%             \addplot[headval] table[x=chi, y expr={\thisrow{head}/\valatchiref}, discard if not={source}{\plainname}, col sep=tab] {\hybridfile};
            \addplot[totval] table[x=chi, y expr={\thisrow{result}/\valatchiref}, discard if not={source}{\plainname}, col sep=tab] {\hybridfile};

%             \legend{ref.,total,head,tail};
        \end{semilogxaxis}
        \begin{loglogaxis}[
            shift={(0,-(\hybridplotheight + \Eplotsep))},
            xmin=\chimin, xmax=\chimax,
            ymin=\errmin,
            width=\hybridplotwidth, height=\hybridplotheight, scale only axis,
            xlabel={$\chi$}, xticklabel pos=bottom,
            ylabel={$\delta\prettyname$}
            ]
            \addplot[tailval] table[x=chi, y=dtail, discard if not={source}{\plainname}, col sep=tab] {\hybridfile};
            \addplot[headval] table[x=chi, y=dhead, discard if not={source}{\plainname}, col sep=tab] {\hybridfile};
            \addplot[totval] table[x=chi, y=dresult, discard if not={source}{\plainname}, col sep=tab] {\hybridfile};
        \end{loglogaxis}

        \useasboundingbox (0,-1.2\hybridplotheight) + (0,-.8cm); % Force larger bounding box
    \end{tikzpicture}}

% Tiny little diagrams for integral subscripts

\newcommand{\scriptdiagrscale}{0.3}
\newlength{\scriptdiagrleglen}\setlength{\scriptdiagrleglen}{0.2cm}
\newlength{\scriptdiagrleg}\setlength{\scriptdiagrleg}{\scriptdiagrleglen}\addtolength{\scriptdiagrleg}{\bubblesize}
\newcommand{\tad}{%
    {\tikzset{external/export next=false}\tikz[scale=\scriptdiagrscale, transform shape]{%
        \tikzset{pion/.style={ thin }}
        \path[use as bounding box] (-1.5\scriptdiagrleglen,0) rectangle (+1.5\scriptdiagrleglen,.6\tadpolesize);
        \draw[pion] (-1.5\scriptdiagrleglen,0) -- (0,0) pic {tadpole} -- (+1.5\scriptdiagrleglen,0);}}}
\newcommand{\bub}{%
    {\tikzset{external/export next=false}\tikz[scale=\scriptdiagrscale, transform shape]{%
        \tikzset{pion/.style={ thin }}
        \draw (0,0) pic {bubble};
        % This (0,0) +  avoids coordinate shenanigans in tikzpicutres
        \draw[pion] (0,0) + (-\bubblesize,0) -- +(-\scriptdiagrleg,0);
        \draw[pion] (0,0) + (+\bubblesize,0) -- +(+\scriptdiagrleg,0);
        }}}
\newcommand{\sunset}{%
    {\tikzset{external/export next=false}\tikz[scale=\scriptdiagrscale, transform shape]{%
        \tikzset{pion/.style={ thin }}
        \renewcommand{\bubbleaspect}{1}
        \draw[pion] (0,0) -- (-\bubblesize,0) -- +(-\scriptdiagrleglen,0);
        \draw[pion] (0,0) -- (+\bubblesize,0) -- +(+\scriptdiagrleglen,0);
        \draw (0,0) pic {bubble};
        }}}
\newcommand{\bball}{%
    {\tikzset{external/export next=false}\tikz[scale=\scriptdiagrscale, transform shape]{%
        \tikzset{pion/.style={ thin }}
        \renewcommand{\bubbleaspect}{1}
        \draw[pion] (0,0) + (-\bubblesize,0) -- +(-\scriptdiagrleg,0);
        \draw[pion] (0,0) + (+\bubblesize,0) -- +(+\scriptdiagrleg,0);
        \draw (0,0) pic {bubble} pic[yscale=.4] {bubble};
        }}}
\newcommand{\fatcat}{%
    {\tikzset{external/export next=false}\tikz[scale=\scriptdiagrscale, transform shape]{%
        \tikzset{pion/.style={ thin }}
        \renewcommand{\bubbleaspect}{1}
        \draw[pion] (0,0) + (-\bubblesize,0) -- +(-\scriptdiagrleg,0);
        \draw[pion] (0,0) + (+\bubblesize,0) -- +(+\scriptdiagrleg,0);
        \draw (0,0) pic {bubble=40};
        \draw (node40) pic[rotate=-40, scale=.6] {bubble};
        }}}
\newcommand{\catseye}{%
    {\tikzset{external/export next=false}\tikz[scale=\scriptdiagrscale, transform shape]{%
        \tikzset{pion/.style={ thin }}
        \renewcommand{\bubbleaspect}{1}
        \draw[pion] (0,0) + (-\bubblesize,0) -- +(-\scriptdiagrleg,0);
        \draw[pion] (0,0) + (+\bubblesize,0) -- +(+\scriptdiagrleg,0);
        \draw (0,0) pic {bubble} pic[xscale=.4] {bubble};
        }}}

\newcommand{\catseyeplus}{%
    {\tikzset{external/export next=false}\tikz[scale=\scriptdiagrscale, transform shape]{%
        \tikzset{pion/.style={ thin }}
        \renewcommand{\bubbleaspect}{1}
        \draw[pion] (0,0) + (-\bubblesize,0) -- +(-\scriptdiagrleg,0);
        \draw[pion] (0,0) + (+\bubblesize,0) -- +(+\scriptdiagrleg,0);
        \draw (0,0) pic {bubble={60,120,270}};
        \draw[pion] (node60)  -- (node270);
        \draw[pion] (node120) -- (node270);
        }}}
\newcommand{\catseyeplusplus}{%
    {\tikzset{external/export next=false}\tikz[scale=\scriptdiagrscale, transform shape]{%
        \tikzset{pion/.style={ thin }}
        \renewcommand{\bubbleaspect}{1}
        \draw[pion] (0,0) + (-\bubblesize,0) -- +(-\scriptdiagrleg,0);
        \draw[pion] (0,0) + (+\bubblesize,0) -- +(+\scriptdiagrleg,0);
        \draw (0,0) pic {bubble={60,120,240,300}};
        \draw[pion] (node60)  to[bend right=30] (node300);
        \draw[pion] (node240) to[bend right=30] (node120);
        }}}


\newcommand{\hypergeometric}[5][]{{}_{2\hspace{-1pt}}\mathit{F}_{\hspace{-1pt}1}#1(\genfrac{}{}{0pt}{}{#2,#3}{#4}#1|#5#1)}
\DeclareRobustCommand{\map}[2]{\mu_{\renewcommand{\scriptdiagrscale}{.2}#1#2}}
\newcommand{\Jbub}[1]{J_\bub^{(#1)}}

\renewcommand{\Re}{\mathrm{Re}}
\renewcommand{\Im}{\mathrm{Im}}
\renewcommand{\d}{\mathrm{d}}
\newcommand{\rat}{\mathrm{rat}}
\newcommand{\Reg}{\mathrm{reg}}
\newcommand{\Div}{\mathrm{div}}
\newcommand{\I}{\mathcal I}
\renewcommand{\IJ}[1]{\I_J^{(#1)}}
\newcommand{\IE}[1]{\I_E^{(#1)}}
\newcommand{\Irat}[1]{\I_\rat^{(#1)}}
\newcommand{\Idiv}{\I_\Div}

% The captions are included in generate_plots.tex so that the size of the figures can be adjusted.
% However, this means that the same caption would have to be maintained both there and where the figure is included in the main text.
% The solution is to break them out into this file.

\newcommand{\figellipticcaption}{%
    \caption{The integrals $E_1$, $E_2$ and $E_3$ (left to right), shown as functions of the nome $q$ over the unit disk in the complex plane.
    For each function, its value at $t=0$ (limit as $q\to1$) has been subtracted to make the features more visible.
    The hue represents the complex argument (with \textbf{\color{red}red} and \textbf{\color{cyan}cyan} for positive and negative real values, respectively), while the shaded contours represent the magnitude logarithmically: with each successive contour, $|t|$ gets a factor of $2$ larger (brighter side) or smaller (darker side).
    Each $t\neq 0$ corresponds to infinitely many $q$ via \cref{eq:q-to-t}, but the wedge-shaped region outlined in white corresponds to a single copy of the complex $t$-plane, with $t=0$ at its rightmost extremity; this is explored further in \cref{sec:t-tau-q}.
    Note the branch cut on $q\in[0,1]$, which corresponds to two copies of $t\in[16,\infty)$ (see \cref{fig:q}).
    }}


\newcommand{\figtaucaption}{%
    \caption{\textbf{Left}: the mapping $\tau\mapsto t$ visualized in the complex plane, covering one period in the real direction.
%     The hue represents the complex argument (with \textbf{\color{red}red} and \textbf{\color{cyan}cyan} for positive and negative real values, respectively), while the shaded contours represent the magnitude logarithmically: with each successive contour, $|t|$ gets a factor of $2$ larger (brighter side) or smaller (darker side).
    The complex values are represented as in \cref{fig:elliptic}.
    The $\times$-shaped contour intersections are saddle points at $t=4$ (for instance at \mbox{$\tau=\pm\frac14+\frac{i\sqrt3}{12}$}) and $t=16$ (for instance at \mbox{$\tau=\pm\frac12+\frac{i\sqrt3}{6}$}).
    \textbf{Right}: a sketch (to scale) of the $\Re(\tau)>0$ half of the same plot, showing a selection of the infinitely many $\tau$-plane images of the real $t$-axis, colored according to which kinematic region they reproduce: spacelike subthreshold (\mbox{$t<0$}, \textbf{\color{spacelike}blue}), timelike subthreshold (\mbox{$0<t<4$}, \textbf{\color{subthr}orange}), two-pion (\mbox{$4<t<16$}, \textbf{\color{twopi}green}) and multi-pion (\mbox{$16<t$}, \textbf{\color{fourpi}violet}).
    The shaded area (if extended upward to infinity) contains all $t$ with $\Im(t)>0$ exactly once, with contours indicating (bottom to top) $\Im(t)=0.001$, $0.1$, $0.5$, and $1$; its boundary (bold) contains all real $t$ exactly once, ordered counterlockwise.}}

\newcommand{\figqcaption}{%
    \caption{Like \cref{fig:tau}, but showing $q\mapsto t$ over the unit disk.
            The lines in the sketch have been extended onto the $\Re(\tau)<0$ part (lower half of the disk), and thereby produce the outline used in \cref{fig:elliptic}.
            Note how inside the shaded region, $q$ approaches $1$ extremely slowly as $|t|$ approaches $0$, ensuring good convergence of our sums almost everywhere.
        }}

\newcommand{\figtsuncaption}{%
    \caption[]{Like \cref{fig:tau}, but showing {\tikzexternaldisable $t_\sunset\mapsto t$} as defined in \cref{eq:t-to-tau}.
        In the sketch on the right, the shaded area is $\Omega_>$, and the two inverse images of the real $t$-line are shown with bold colored lines as in \cref{fig:tau}.
        The dotted rectangle indicates the region shown in \cref{fig:rho24} (bottom row) and \cref{fig:rhoF}.
        The thin gray lines are the same contours of constant $\Im(t)$ as in \cref{fig:tau}, and the dashed ones indicate their images under the involutory mapping $z\mapsto (9z-9)/(z-9)$, which is how they appear in $\varpi_r$.}}


\newcommand{\figrhotwofourcaption}{%
    \caption{
        The construction of $\rho_2(z)$ (top row) and $\rho_4(z)$ (bottom row).
        \textbf{Left}: The principal branch of $\sqrt{(z-4)(z-16)}$ and $\sqrt[4]{z^3 - 9z^2 + 3z - 3}$, respectively, visualized in the complex plane using colors and contours like in \cref{fig:tau}.
        \textbf{Middle}: The corrected functions $\rho_2(z)$ and $\rho_4(z)$, respectively.
        \textbf{Right}: A schematic picture of their relation.
        Colored lines indicate where the argument of the root [$(z-4)(z-16)$ and $z^3 - 9z^2 + 3z - 3$, respectively] is real: \textbf{\color{plotV} red} for positive values and \textbf{\color{plotII} cyan} for negative.
        Dotted lines do not cause branch cuts.
        Solid lines cause cuts that are removed by analytically continuing to different branches, starting on the principal branch in the region marked by $\star$ and picking up roots of unity ($-1$ and $+i$, respectively) when crossing a line in the direction indicated by the arrows.
        Zigzag lines are left as branch cuts also in $\rho_2(z)$ or $\rho_4(z)$.
        The dotted circle indicates the boundary between $\Omega_<$ and $\Omega_>$ (compare \cref{fig:tsun}).
%             It can be noted that a version with cuts only along the dotted red lines is possible by assigning sheets differently, but that is not suitable for $t\mapsto\tau$.
        }}

\newcommand{\figrhoFcaption}{%
    \caption[The function $\rho_F(z)-F_1(1)$, visualized over the same region as $\rho_4(z)$ in \cref{fig:rho24}.]{
        \textbf{Top}:
        The function $\rho_F(z)-F_1(1)$, visualized over the same region as $\rho_4(z)$ in \cref{fig:rho24}.
        Subtracting the constant $F_1(1)$ makes the features more visible.
        The upper half-plane shows $\rho_F(z)$ as used for the
        % $\map{t}{\tau}$
        $t\mapsto\tau$
        relation, and the lower shows the uncorrected $F_1\big(1728/j_\sunset(z)\big)$ using the principal branch of the hypergeometric function.
        Both are symmetric (up to phases) across the real axis.
        The insets show the region around $z=9$, which for the uncorrected function contains the only discontinuity outside $\Omega_<$, making its removal the most important part of the construction.
%         The discontinuity in the magnified region around $z=9$ has a particularly devastating effect on
%         $\map{t}{\tau}$ since it is contained in $\Omega_>$, not $\Omega_<$,
%         % $t\mapsto\tau$,
%         % since the image of the real line under the mapping
%         % $\map{t}{t_\sunset}$
%         % $t\mapsto t_\sunset$
%         % passes through it in a way that
%         and unless remedied, it
%         leads to very problematic interplay between $\varpi_c$ and $\varpi_r$.
        \textbf{Bottom}:
        A schematic picture of the construction of the correct $\rho_F(z)$,
        using the same line conventions as in \cref{fig:rho24}: dotted for no cut, solid for removed cut, zigzag for remaining cut.
        Negative real $z$ values only produce a cut for $F_2(z)$ and therefore do not affect the principal branch ($\star$), which is purely $F_1(z)$.
        Positive real $z$ values only produce a cut when \mbox{$|1728/j_\sunset(z)|\geq 1$}, which holds within the \textbf{\color{realextra} orange} area.
        The monodromies described in the text are used when crossing the solid lines to keep $\rho_F(z)$ continuous: $C_0$ ($\rightarrow$) or $C_1$ ($\leftrightarrow$) is applied when going in the direction of the arrows (recall that $C_1$ is involutory, hence the bidirectional arrow).
        }}

\newcommand{\figrhoFconstrcaption}{%
    \caption[ A schematic version of \cref{fig:rhoF}, showing the construction of the correct $\rho_F(z)$.]{
        A schematic version of \cref{fig:rhoF}, showing the construction of the correct $\rho_F(z)$.
        The lines indicate where $1728/j_\sunset(z)$ (the argument of the hypergeometric functions) is real, using the same conventions as in \cref{fig:rho24}: dotted for no cut, solid for removed cut, zigzag for remaining cut.
        Negative real values only produce a cut for $F_2(z)$ and therefore do not affect the principal branch ($\star$), which is purely $F_1(z)$.
        Positive real values only produce a cut when $|1728/j_\sunset(z)|\geq 1$, which holds within the shaded orange area.
        The monodromies described in the text are used when crossing the solid lines to keep $\rho_F(z)$ continuous: $C_0$ ($\rightarrow$) or $C_1$ ($\leftrightarrow$) is applied when going in the direction of the arrows (recall that $C_1$ is its own inverse, hence the bidirectional arrow).
        }}

\newcommand{\fignumericexamplecaption}{%
    \caption{
        Example of the process of computing $\bar E_5$ and $\bar E_6$ at $t_\ex=4.2+i\epsilon$ following \cref{eq:Ebar-detailed}, with the final summation following \cref{eq:Ebar-schematic}.
        The integration contours are shown in the $\tau$ and $\xi$ planes, which are represented as in \cref{fig:tau}.
        Filled triangles indicate a complex integration contour, and empty ones a real integral toward infinity (using $\theta$ for that along the imaginary $\tau$ axis).
        The shaded area indicates the ``too close to the cut'' zone described in the definition of $\mathcal C_\xi(\beta)$.
        For integrals using \cref{eq:hybrid-integral}, the parts from below ($<$) and above ($>$) the cutoff $\chi$ are indicated.
        % All numbers are truncated; see \cref{tab:numeric-example} for the full numbers.
        All numbers are much more precise than shown here, with the largest errors [those of $\IE{n}(\beta_\star)$] being on the order of $10^{-6}$.
        % Note the presence of large cancellations, with several intermediate quantities being several orders of magnitude larger than the end result; this is even more striking when $t\in[0,4)$, where the large imaginary parts of the contours cancel to give a real result.
        % (Of course, if the equivalent method described at the end of \cref{sec:spacelike} is used, everything is manifestly real).
        }}

\newcommand{\fighybridcaption}{%
    \caption{
        The four usages of \cref{eq:hybrid}, evaluated at $\beta=\beta_\star$ and shown as functions of the cutoff $\chi$.
%           In each case, the top panel shows the total value (\linesample{totval}) along with that of the ``head'' (\linesample{headval}), i.e., $\xi<\chi$ numerically integrated, and the ``tail'' (\linesample{tailval}), i.e., $\xi>\chi$ using a series expansion.
        The bottom panel shows the absolute error estimate of each part.
        For easier comparison, all values (but not the errors) have been divided by the corresponding total value of the function evaluating using our preferred $\chi$.
        This $\chi$, which is what we use elsewhere, is marked with a dot.
        }}

\newcommand{\figmonodromycaption}{%
    \caption{The contours used in the monodromy construction, with cuts drawn as zigzag lines for the different functions involved.
        \textbf{Left}: $F_1(z)$, cut on $[1,\infty)$.
        \textbf{Middle}: $F_2(z)$, cut on $(-\infty,0]$.
        \textbf{Right}: $\hypergeometric abc{\frac1z}$, showing how the combination of $C_0$ and $C_1$ results in a contour that encircles the cut without crossing it.
        }}

\newcommand{\figJbubcaption}{%
    \caption[The function Jbub.]{
        The functions $\Jbub{n}(t)$, normalized so that the magnitude is roughly constant in $n$.
        Real parts are shown as solid lines, and imaginary parts as dashed lines; colors indicate different $n$ as per the plot legend.
        Below the 2-pion threshold ($t=4$), the imaginary parts are identically zero and are not shown.
        Above threshold, $t$ has been given a small positive imaginary part.
        }}

\newcommand{\figmasterscaption}{%
    \caption[The six elliptic master integrals plotted as functions of $t$.]{
        The six elliptic master integrals plotted as functions of $t$, showing the real (\Resample) and imaginary (\Imsample) parts, as well as the fifth-order series expansion around $t=0$ (\seriessample).
%         The two- and four-pion thresholds are drawn as vertical dotted lines.
%         Below each plot is shown the relative difference between \texttt{pySecDec} (lines with error bands) and our implementation, whose error is not visible at the scale of the plots.
%         Also shown is the relative difference between our implementation and the series expansion.
%         Reinterpreting the below-threshold imaginary part as errors (like in \cref{fig:abserr}) gives error bands much taller than the plots.
        At the scale of the plots, the difference between our implementation and \texttt{pySecDec} is not visible (but see \cref{fig:reldiff}), nor are the error bands.
        However, the absolute uncertainty of each line is shown on a logarithmic scale below each plot, using \texttt{pySecDec}'s uncertainty; the uncertainty using our method is shown as a shaded area.
        For the series expansion, the truncation error is estimated as the magnitude of the last non-truncated term.
        }}

\newcommand{\figreldiffcaption}{%
    \caption[Symmetric log plots of the relative difference between the three different implementations of the master integrals]{
        Symmetric log plots of the relative difference between our results and those of \texttt{pySecDec} and the series expansion, shown similarly to \cref{fig:masters}.
        By relative difference, we mean $[x - x_\text{ref}]/x_\text{ref}$ where $x$ is some quantity and $x_\text{ref}$ is the same quantity as computed using our method.
        Below threshold, the imaginary part is known to be zero and the relative difference makes little sense, so in the parts of the plots that are shaded ({\setbox0=\hbox{Xg}\tikzset{external/export next=false}\tikz[baseline=\dp0]{\draw[gray, pattern=north east lines, pattern color=gray] (0,0) rectangle (\dimexpr\ht0+\dp0\relax,\dimexpr\ht0+\dp0\relax);}}) we display the actual quantities $x$ rather than the relative difference.
        Our values (whose relative difference is zero by construction) and their error bands are shown in gray.
        }}

\newcommand{\figtimecaption}{%
    \caption[Time consumption, measured in seconds per master integral evaluation, needed by the different evaluation methods for different $t$.]{
        Time, measured in seconds per sample, needed to obtain the data behind \cref{fig:masters,fig:reldiff}.
        Solid lines indicate our implementation, and dashed lines indicate \texttt{pySecDec}; colors represent different masters as per the plot legend.
        Note that our implementation is single-threaded, whereas \texttt{pySecDec} made full use of the 16-core processor on which this data was collected.
        }}

\newcommand{\figPiEcaption}{%
    \caption[$\bar\Pi_E$]{
        The real (solid) and imaginary (dashed) parts of $\bar\Pi_E(t)$, and its fifth-order series expansion around $t=0$ (dotted), which has radius of convergence $4$.
        The error bands are not visible on the scale of the plot, but the absolute uncertainties are shown on a logarithmic scale below the main panel.
        Also shown are the contributions from individual master integrals (i.e., the term in \cref{eq:PiE} consisting of that integral times a Laurent polynomial of $t$), as well as the $E_i$-independent remainder.
        The imaginary parts are only shown above their respective thresholds.
        The error bands are too small to be seen at the scale of the plot, but the absolute uncertainties are shown on a logarithmic scale below the main plot.
        }}

\newcommand{\figPiJcaption}{%
    \caption[$\bar\Pi_J$]{
        $\bar\Pi_J(t)$, shown similarly to $\bar\Pi_E(t)$ in \cref{fig:PiE}.
        Here, the contributions are separated based on proportionality to different products of $\Jbub{n}$.
        Also shown is the purely real $\bar\Pi_\zeta(t)$ (dotted).
        The uncertainties are minuscule and are not shown.
        }}


\newcommand{\figtauerrcaption}{%
    \caption[]{
        The numerical error $\delta\tau\coloneq\map{t}{\tau}\big[\map{\tau}{t}(\tau)\big]-\tau$ obtained when mapping $\tau$ to $t$ and back, shown in the $\tau$ plane similarly to \cref{fig:tau}.
        Since $\map{\tau}{t}$ is exact to within negligible arithmetic noise
        % (128-bit floating point numbers are used here, or $\gtrsim35$ digits of precision),
        this effectively shows the precision of $\map{t}{\tau}$, which is clearly dominated by a smooth systematic error rather than noise.
        The left panel illustrates the phase of $\delta\tau$ as in figs.~\ref{fig:tau},~\ref{fig:q}, etc., whereas the right shows only the absolute value to more clearly demonstrate its size.
        The jump in error marked by the arrow corresponds to the argument of $\varpi_r$ crossing the branch cut near $z=9$ that is shown magnified in \cref{fig:rhoF}; although the discontinuity of $\rho_F$ has been removed using the monodromy, the introduction of $F_2$ of course changes its uncertainty.
        Note also the conspicuous row of four zeroes (dark spots in the right panel).
        These lie along the curve $|t_\sunset-9|=\sqrt{72}$, on which \cref{eq:t-to-tau} reduces to $i\varpi_c(t_\sunset)/\varpi_r(t_\sunset^*)$; evidently, this allows the errors to more or less cancel.}}



\newcommand{\plotdir}{../plots}

\usetikzlibrary{external}
\tikzsetexternalprefix{tikz/}
\tikzexternalize

\begin{document}

%%% For elliptic.tex %%%

\begin{figure}[tp]
    \centering
%     \tikzsetnextfilename{fig-elliptic-w1}
%     \qplot{w1.png}
    \tikzsetnextfilename{fig-elliptic-E1}
    \qplot{E1q.png}
    \tikzsetnextfilename{fig-elliptic-E2}
    \qplot{E2q.png}
    \tikzsetnextfilename{fig-elliptic-E3}
    \qplot{E3q.png}
    \figellipticcaption
    \label{fig:elliptic}
\end{figure}

\clearpage

%%% For t-tau-q.tex %%%

\begin{figure}[p]
    \begin{subfigure}[t]{.6\columnwidth}
        \hspace{-.7cm}
        \tikzsetnextfilename{fig-tau-left}
        \begin{tikzpicture}
            \begin{axis}[
                width=\columnwidth, height=.5\columnwidth, scale only axis, axis on top,
                enlargelimits=false,
                axis on top,
                xlabel={$\Re(\tau)$},
                ylabel={$\Im(\tau)$},
                xtick={-.5,-.25,0,+.25,+.5},
                xticklabels={$-\tfrac12$, $-\tfrac14$,$0$,$\tfrac14$,$\tfrac12$},
                ytick={0,\suprad/2,\suprad,.5},
                yticklabels={$0$,$\tfrac{\sqrt3}{12}$,$\tfrac{\sqrt3}{6}$,$\tfrac12$}
                ]
                \addplot graphics [
                    xmin=-.5, xmax=.5,
                    ymin=0, ymax=.5,
                ] {\plotdir/tau.png};
            \end{axis}
        \end{tikzpicture}
    \end{subfigure}
    \hspace{.7cm}
    \begin{subfigure}[t]{.3\columnwidth}
        \tikzsetnextfilename{fig-tau-right}
        \begin{tikzpicture}
            \begin{axis}[
                width=\columnwidth, height=\columnwidth, scale only axis, clip=false,
                xmin=0, xmax=.5,
                ymin=0, ymax=.5,
                xlabel={$\vphantom{\Re(\tau)}$},
                xtick={0,.111,.1667,.25,.333,.5},
                xticklabels={$0$,$\tfrac19$,$\tfrac16$,$\tfrac14$, $\tfrac13$, $\tfrac12$},
                ytick={0,\suprad/2,\suprad,.5},
                yticklabels={$0$,$\tfrac{\sqrt3}{12}$,$\tfrac{\sqrt3}{6}$,$\tfrac12$},
                yticklabel pos=right
                ]

                % Shaded principal inverse image of the upper half-plane
                \fill[shaded] (0,.5) -- (0,0)
                    arc[start angle=180, end angle=60, radius=\subrad]
                    arc[start angle=150, end angle=90, radius=\suprad]
                    -- (.5,.5);
                \draw[gray,densely dashed] (.25,0) -- (.25,{\suprad/2});

                % A selection of inverse images of segments of the real line
                \draw[spacelike] (0,0) arc[start angle=180, end angle=0, radius=\subrad/2];
                \draw[spacelike] (0,0) arc[start angle=180, end angle=0, radius=\subrad/4];
                \draw[spacelike] (0,0) arc[start angle=180, end angle=0, radius=\subrad/6];
                \draw[spacelike] (0,0) arc[start angle=180, end angle=0, radius=\subrad/8];
                \draw[spacelike] (0,0) arc[start angle=180, end angle=0, radius=\subrad/10];
                \draw[subthr] (0,0) arc[start angle=180, end angle=0, radius=\subrad/1];
                \draw[subthr] (0,0) arc[start angle=180, end angle=0, radius=\subrad/3];
                \draw[subthr] (0,0) arc[start angle=180, end angle=0, radius=\subrad/5];
                \draw[subthr] (0,0) arc[start angle=180, end angle=0, radius=\subrad/7];
                \draw[subthr] (0,0) arc[start angle=180, end angle=0, radius=\subrad/9];
                \draw[subthr] (0,0) arc[start angle=180, end angle=0, radius=\subrad/11];
                \draw[thick, fourpi] (+.5,0) -- (+.5,.5);
                \draw[twopi] (.5,\suprad) arc[start angle=90, end angle=180, radius=\suprad];
                \draw[spacelike] (.5-\subrad,0) arc[start angle=180, end angle=0, radius=\subrad/2];
                \draw[spacelike] (.5-\subrad/3,0) arc[start angle=180, end angle=0, radius=\subrad/6];
                \draw[spacelike] (.5-\subrad/5,0) arc[start angle=180, end angle=0, radius=\subrad/10];
                \draw[fourpi] (\subrad,0) arc[start angle=180, end angle=0, radius=\subrad/4];
                \draw[fourpi] (.5-\subrad/2,0) arc[start angle=180, end angle=0, radius=\subrad/4];
                \draw[fourpi] (.5-\subrad/4,0) arc[start angle=180, end angle=0, radius=\subrad/8];
                \draw[fourpi] (.5-\subrad/6,0) arc[start angle=180, end angle=0, radius=\subrad/12];
                \draw[spacelike] (1/4,0) arc[start angle=180, end angle=0, radius=\subrad/4];

                % The principal inverse image of the real line
                \draw[realline, spacelike] (0,.5) -- (0,0);
                \draw[realline, subthr] (0,0)
                    arc[start angle=180, end angle=60, radius=\subrad]
                    node[inner sep=0, anchor=center, pin=0:{\tiny\color{black} $t=4$}] {};
                \draw[realline, twopi] (.5,\suprad)
                    node[inner sep=0, anchor=center, pin=240:{\tiny\color{black} $t=16$}] {}
                    arc[start angle=90, end angle=150, radius=\suprad];
                \draw[realline, fourpi] (.5,\suprad) -- (.5,.5);

                % Contours at fixed epsilon
                \begin{scope}
                    \clip (0,0) rectangle (.5,.5);
                    \addplot[very thin, gray, mark=none] table [x={Retau},y={Imtau}] {\plotdir/image_Imt1.0.dat} -- cycle;
                    \addplot[very thin, gray, mark=none] table [x={Retau},y={Imtau}] {\plotdir/image_Imt0.5.dat} -- cycle;
                    \addplot[very thin, gray, mark=none] table [x={Retau},y={Imtau}] {\plotdir/image_Imt0.1.dat} -- cycle;
                    \addplot[very thin, gray, mark=none] table [x={Retau},y={Imtau}] {\plotdir/image_Imt0.001.dat} -- cycle;
                \end{scope}

                % Annotations
                \node[anchor=west,inner sep=0] at (.03,.45) {\tiny $\uparrow -\infty$};
                \node[anchor=west,inner sep=0] at (.03,.40) {\tiny $\downarrow t=0$};
                \node[anchor=east,inner sep=0] at (.48,.45) {\tiny $+\infty\uparrow $};
            \end{axis}
        \end{tikzpicture}
    \end{subfigure}
    \figtaucaption
    \label{fig:tau}
\end{figure}


\begin{figure}[p]
    \centering
    \tikzsetnextfilename{fig-q-left}
    \begin{subfigure}[t]{.4\columnwidth}
        \begin{tikzpicture}
            \begin{axis}[
                width=\columnwidth, height=\columnwidth, scale only axis, axis on top,
                xmin=-1, xmax=1,
                ymin=-1, ymax=1,
                ticks=none, axis lines=none
                ]
                \addplot graphics [
                    xmin=-1, xmax=1,
                    ymin=-1, ymax=1
                ] {\plotdir/q.png};
            \end{axis}
        \end{tikzpicture}
    \end{subfigure}
    \hspace{1cm}
    \begin{subfigure}[t]{.4\columnwidth}
        \tikzsetnextfilename{fig-q-right}
        \begin{tikzpicture}
            \begin{axis}[
                width=\columnwidth, height=\columnwidth, scale only axis,
                xmin=-1, xmax=1,
                ymin=-1, ymax=1,
                ticks=none, axis lines=none
                ]
                \clip (0,0) circle[radius=1];

                \addqarcs[spacelike,domain=0:180]{\subrad/2}{\subrad/2}
                \addqarcs[spacelike,domain=0:180]{\subrad/4}{\subrad/4}
                \addqarcs[spacelike,domain=0:180]{\subrad/6}{\subrad/6}
                \addqarcs[spacelike,domain=0:180]{\subrad/8}{\subrad/8}
                \addqarcs[spacelike,domain=0:180]{\subrad/10}{\subrad/10}
                \addqarcs[subthr,domain=0:180]{\subrad/1}{\subrad/1}
                \addqarcs[subthr,domain=0:180]{\subrad/3}{\subrad/3}
                \addqarcs[subthr,domain=0:180]{\subrad/5}{\subrad/5}
                \addqarcs[subthr,domain=0:180]{\subrad/7}{\subrad/7}
                \addqarcs[subthr,domain=0:180]{\subrad/9}{\subrad/9}
                \addqarcs[subthr,domain=0:180]{\subrad/11}{\subrad/11}
                \addqarcs[twopi,domain=0:180,samples=32]{.5}{\suprad}
                \draw[fourpi] (-1,0) -- (-.1,0);
                \addqarcs[spacelike,domain=0:180]{.5-\subrad/2}{\subrad/2}
                \addqarcs[spacelike,domain=0:180]{.5-\subrad/6}{\subrad/6}
                \addqarcs[spacelike,domain=0:180]{.5-\subrad/10}{\subrad/10}
                \addqarcs[fourpi,domain=0:180]{\subrad*5/4}{\subrad/4}
                \addqarcs[fourpi,domain=0:180]{.5-\subrad/4}{\subrad/4}
                \addqarcs[fourpi,domain=0:180]{.5-\subrad/8}{\subrad/8}
                \addqarcs[fourpi,domain=0:180]{.5-\subrad/12}{\subrad/12}
                \addqarcs[spacelike,domain=0:180]{\subrad*7/4}{\subrad/4}

                \draw[gray,thick] (0,0) circle[radius=1];
                \draw[gray] (0,-1) -- (0,+1);

                \addplot[name path=principal fourpi, realline, fourpi, domain=-1:0, samples=2] ({x*exp(-2*pi*\suprad)},0);
                \addqarc[name path=principal subthr, realline, subthr, domain=60:180]{\subrad}{\subrad}
                \addqarc[name path=principal twopi, realline, twopi, domain=90:150]{.5}{\suprad}
                \addplot[name path=principal spacelike, realline, spacelike, domain=0:1, samples=2] (x,0);

                \addplot[shaded] fill between [of=principal fourpi and principal twopi];
                \addplot[shaded] fill between [of=principal subthr and principal spacelike];

                \addplot[very thin, gray, mark=none] table [x={Req},y={Imq}] {\plotdir/image_Imt1.0.dat} -- cycle;
                \addplot[very thin, gray, mark=none] table [x={Req},y={Imq}] {\plotdir/image_Imt0.5.dat} -- cycle;
                \addplot[very thin, gray, mark=none] table [x={Req},y={Imq}] {\plotdir/image_Imt0.1.dat} -- cycle;
                \addplot[very thin, gray, mark=none] table [x={Req},y={Imq}] {\plotdir/image_Imt0.001.dat} -- cycle;

                \node[anchor=center, inner sep=0, pin=-45:{\tiny$\pm\infty$}] at (0,0) {};
                \node[anchor=center, inner sep=0, pin=160:{\tiny$t=16$}] at ({-exp(-2*pi*\suprad)},0) {};
                \node[anchor=center, inner sep=0, pin=180:{\tiny$t=4$}] at (0,{exp(-pi*\suprad}) {};
            \end{axis}
        \end{tikzpicture}
    \end{subfigure}
    \figqcaption
    \label{fig:q}
\end{figure}

%%% For tsun.tex %%%

\begin{figure}[p]
    \begin{subfigure}[t]{.45\columnwidth}
        \hspace{-.7cm}
        \tikzsetnextfilename{fig-tsun-left}
        \begin{tikzpicture}
            \begin{axis}[
                width=\columnwidth, height=.417\columnwidth, scale only axis, axis on top,
                xmin=-8, xmax=16,
                ymin=-5, ymax=5,
                enlargelimits=false,
                xlabel={$\Re(t_\sunset)$},
                ylabel={$\Im(t_\sunset)$},
                xtick={-6,-3,0,3,6,9,12,15},
                ytick={-3,0,3},
                ]
                \addplot graphics [
                    xmin=-8, xmax=16,
                    ymin=-5, ymax=5,
                ] {\plotdir/tsun.png};
            \end{axis}
        \end{tikzpicture}
    \end{subfigure}
    \hspace{.7cm}
    \begin{subfigure}[t]{.45\columnwidth}
        \tikzsetnextfilename{fig-tsun-right}
        \begin{tikzpicture}[pin distance=2mm]
            \begin{axis}[
                width=\columnwidth, height=.417\columnwidth, scale only axis,
                xmin=-8, xmax=16,
                ymin=-5, ymax=5,
                xlabel={$\vphantom{\Re(\tau)}$},
                xtick={0,.111,.1667,.25,.333,.5},
                xtick={-6,-3,0,3,6,9,12,15},
                ytick={-3,0,3},
                yticklabels={},
                ]

                % Shaded principal inverse image of the upper half-plane
                \fill[shaded, even odd rule] (-8,-5) rectangle (16,5) (0,0) circle[radius=3];
                % Footprint of other plots
                \draw[gray,densely dotted] (-3,-3) rectangle (11,3);

                % The two inverse images of the real line
                \draw[realline, spacelike] (0,0) -- (1,0);
                \draw[realline, spacelike] (9,0) -- (20,0);
                \draw[realline, subthr] (-3,0) -- (0,0)
                    node[inner sep=0, anchor=center, pin=90:{\tiny\color{black} $t=0$}] {};
                \draw[realline, subthr] (-9,0) -- (-3,0);
                \draw[realline, twopi] (3,0) arc[start angle=0, end angle=180, radius=3]
                    node[inner sep=0, anchor=center, pin=240:{\tiny\color{black} $t=4$}] {};
                \draw[realline, twopi] (-3,0) arc[start angle=180, end angle=360, radius=3]
                    node[inner sep=0, anchor=center, pin=60:{\tiny\color{black} $t=16$}] {};
                \draw[realline, fourpi] (3,0) -- (1,0)
                    node[inner sep=0, anchor=center, pin=270:{\tiny\color{black} $\pm\infty$}] {};
                \draw[realline, fourpi] (3,0) -- (9,0)
                    node[inner sep=0, anchor=center, pin=270:{\tiny\color{black} $\pm\infty$}] {};

%                 Contours at fixed epsilon
                \addplot[tcontour=-5:5] table [x={Rezc},y={Imzc}] {\plotdir/image_Imt1.0.dat} -- (9,0);
                \addplot[tcontour=-5:5] table [x={Rezc},y={Imzc}] {\plotdir/image_Imt0.5.dat} -- (9,0);
                \addplot[tcontour=-5:5] table [x={Rezc},y={Imzc}] {\plotdir/image_Imt0.1.dat} -- (9,0);
                \addplot[tcontour=-5:5] table [x={Rezc},y={Imzc}] {\plotdir/image_Imt0.001.dat} -- (9,0);

                \addplot[tcontour=-5:5, densely dashed] table [x={Rezr},y={Imzr}] {\plotdir/image_Imt1.0.dat};
                \addplot[tcontour=-5:5, densely dashed] table [x={Rezr},y={Imzr}] {\plotdir/image_Imt0.5.dat};
                \addplot[tcontour=-5:5, densely dashed] table [x={Rezr},y={Imzr}] {\plotdir/image_Imt0.1.dat};
                \addplot[tcontour=-5:5, densely dashed] table [x={Rezr},y={Imzr}] {\plotdir/image_Imt0.001.dat};

                % Annotations
                \node[anchor=west,inner sep=0] at (-7.5,.40) {\tiny $\leftarrow t=0$};
                \node[anchor=east,inner sep=0] at (15.5,.40) {\tiny $t=0\rightarrow$};
            \end{axis}
        \end{tikzpicture}
    \end{subfigure}
    \figtsuncaption
    \label{fig:tsun}
\end{figure}

%%% For rho24.tex %%%

\begin{figure}[p!]
    \centering
    \hspace{-.7cm}
    \tikzsetnextfilename{fig-rho2-left}
    \begin{tikzpicture}
        \begin{axis}[
                width=.3\columnwidth, height=.129\columnwidth, scale only axis, axis on top,
                xmin=-4, xmax=24,
                ymin=-6, ymax=6,
                xlabel={$\Re(z)$},
                ylabel={$\Im(z)$},
                ytick={-4,0,4},
                xticklabel pos=top
            ]
            \addplot graphics [
                xmin=-4, xmax=24,
                ymin=-6, ymax=6,
            ] {\plotdir/rho2_before.png};
        \end{axis}
    \end{tikzpicture}
    \hspace{-.5cm}
    \tikzsetnextfilename{fig-rho2-mid}
    \begin{tikzpicture}
        \begin{axis}[
                width=.3\columnwidth, height=.129\columnwidth, scale only axis, axis on top,
                xmin=-4, xmax=24,
                ymin=-6, ymax=6,
                xlabel={$\Re(z)$},
                ylabel={},
                ytick={-4,0,4},
                yticklabels={},
                xticklabel pos=top
            ]
            \addplot graphics [
                xmin=-4, xmax=24,
                ymin=-6, ymax=6,
            ] {\plotdir/rho2_after.png};
        \end{axis}
    \end{tikzpicture}
    \hspace{-.5cm}
    \tikzsetnextfilename{fig-rho2-right}
    \begin{tikzpicture}
        \begin{axis}[
                width=.3\columnwidth, height=.129\columnwidth, scale only axis,
                xmin=-4, xmax=24,
                ymin=-6, ymax=6,
                xlabel={$\Re(z)$},
                ylabel={$\Im(z)$},
                ylabel={},
                ytick={-4,0,4},
                yticklabels={},
                xticklabel pos=top
            ]

            \coordinate (zim1) at (4,0);
            \coordinate (zim2) at (16,0);

            \draw[realneg, branchcut] (4,0) -- (16,0);
            \draw[realneg] (10,6) -- (10,-6);
            \draw[realpos, nocut] (-4,0) -- (4,0) (16,0) -- (24,0);

            \draw[C1] (8,+3) -- (12,+3);
            \draw[C1] (8,-3) -- (12,-3);

            \node at (18,2) {$\star$};
        \end{axis}
    \end{tikzpicture}

    \hspace{-.7cm}
    \tikzsetnextfilename{fig-rho4-left}
    \begin{tikzpicture}
        \begin{axis}[
                width=.3\columnwidth, height=.129\columnwidth, scale only axis, axis on top,
                xmin=-3, xmax=11,
                ymin=-3, ymax=3,
                xlabel={$\Re(z)$},
                ylabel={$\Im(z)$},
            ]
            \addplot graphics [
                xmin=-3, xmax=12,
                ymin=-3, ymax=3,
            ] {\plotdir/rho4_before.png};
        \end{axis}
    \end{tikzpicture}
    \hspace{-.5cm}
    \tikzsetnextfilename{fig-rho4-mid}
    \begin{tikzpicture}
        \begin{axis}[
                width=.3\columnwidth, height=.129\columnwidth, scale only axis, axis on top,
                xmin=-3, xmax=11,
                ymin=-3, ymax=3,
                xlabel={$\Re(z)$},
                ylabel={},
                yticklabels={}
            ]
            \addplot graphics [
                xmin=-3, xmax=11,
                ymin=-3, ymax=3,
            ] {\plotdir/rho4_after.png};
        \end{axis}
    \end{tikzpicture}
    \hspace{-.5cm}
    \tikzsetnextfilename{fig-rho4-right}
    \begin{tikzpicture}
        \begin{axis}[
                width=.3\columnwidth, height=.129\columnwidth, scale only axis,
                xmin=-3, xmax=11,
                ymin=-3, ymax=3,
                xlabel={$\Re(z)$},
                ylabel={$\Im(z)$},
                ylabel={},
                yticklabels={}
            ]

            \coordinate (zim1) at ({3-2*sqrt(2)},0);
            \coordinate (zim2) at ({3+2*sqrt(2)},0);
            \coordinate (root1u) at ({-1/2*4^(2/3)*(1) - 4^(1/3)*(1) + 3},{-1/2*4^(2/3)*(sqrt(3)) - 4^(1/3)*(-sqrt(3))});
            \coordinate (root1l) at ({-1/2*4^(2/3)*(1) - 4^(1/3)*(1) + 3},{1/2*4^(2/3)*(sqrt(3)) + 4^(1/3)*(-sqrt(3))});
            \coordinate (root2) at ({4^(2/3) + 2*4^(1/3) + 3},0);

            \begin{scope}
                \clip (3,3) rectangle (11,-3);
                \addplot[mark=none, realneg] table {\plotdir/im0.dat};
            \end{scope}
            \begin{scope}
                \clip (root1u) -| (3,+3) -- (-3,+3) |- cycle;
                \addplot[mark=none, realpos, nocut] table {\plotdir/im0.dat};
            \end{scope}
            \begin{scope}
                \clip (root1l) -| (3,-3) -- (-3,-3) |- cycle;
                \addplot[mark=none, realpos, nocut] table {\plotdir/im0.dat};
            \end{scope}
            \begin{scope}
                \clip (root1u) -| (3,0) |- (root1l) -| (-3,0) |- cycle;
                \addplot[mark=none, black] table {\plotdir/im0.dat};
                \draw[realneg, branchcut] (root1u) .. controls (.2,0) .. (root1l);
            \end{scope}
            \draw[realneg, branchcut] (-3,0) -- (root2);
            \draw[realpos, nocut] (root2) -- (11,0);

            \draw[C0] (7,+1.0) to[bend right=10] (5,+1.2);
            \draw[C0] (5,-1.2) to[bend right=10] (7,-1.0);

            % Omega boundary
            \draw[gray, densely dotted] (0,0) circle[radius=3];

            \node at (4,-1.5) {$\star$};
        \end{axis}
    \end{tikzpicture}
    \figrhotwofourcaption
    \label{fig:rho24}
\end{figure}

%%% For rhoF.tex %%%

\begin{figure}[p!]
    \def\insetxmin{8.5}
    \def\insetxmax{9.25}
    \def\insetymax{0.25}
    \def\insetXmin{8}
    \def\insetXmax{11}
    \def\insetYmin{2}
    \def\insetYmax{3}

    \centering
    \hspace{-.5cm}
    \tikzsetnextfilename{fig-rhoF-top}
    \begin{tikzpicture}
        \begin{axis}[
                width=.9\columnwidth, height=.3857\columnwidth, scale only axis, axis on top,
                xmin=-3, xmax=11,
                ymin=-3, ymax=3,
                xlabel={$\Re(z)$},
                ylabel={$\Im(z)$},
                xtick={-2,0,2,4,6,8,10},
                xticklabel pos=top
            ]

            % Main plots
            \addplot graphics [
                xmin=-3, xmax=12,
                ymin=-3, ymax=0,
            ]   {\plotdir/rhoF_before.png};
            \addplot graphics [
                xmin=-3, xmax=12,
                ymin=0, ymax=3,
            ] {\plotdir/rhoF_after.png};

            % Insets
            \addplot graphics [
                xmin=\insetXmin, xmax=\insetXmax,
                ymin=+\insetYmin, ymax=+\insetYmax,
            ] {\plotdir/rhoF_after_inset.png};
            \addplot graphics [
                xmin=\insetXmin, xmax=\insetXmax,
                ymin=-\insetYmin, ymax=-\insetYmax,
            ] {\plotdir/rhoF_before_inset.png};

            % Inset lines
            \begin{scope}[gray]
                \draw (\insetxmin,-\insetymax) rectangle (\insetxmax,\insetymax);
                \draw (\insetXmin,+\insetYmax) |- (\insetXmax,+\insetYmin);
                \draw (\insetxmin,+\insetymax) -- (\insetXmin,+\insetYmin);
                \draw (\insetxmax,+\insetymax) -- (\insetXmax,+\insetYmin);
                \draw (\insetXmin,-\insetYmax) |- (\insetXmax,-\insetYmin);
                \draw (\insetxmin,-\insetymax) -- (\insetXmin,-\insetYmin);
                \draw (\insetxmax,-\insetymax) -- (\insetXmax,-\insetYmin);
            \end{scope}

            % Separating line
            \draw (-3,0) -- (11,0);

        \end{axis}
    \end{tikzpicture}

    \hspace{-.5cm}
    \tikzsetnextfilename{fig-rhoF-bot}
    \begin{tikzpicture}
        \begin{axis}[
                width=.9\columnwidth, height=.3857\columnwidth, scale only axis,
                xmin=-3, xmax=11,
                ymin=-3, ymax=3,
                xlabel={$\Re(z)$},
                ylabel={$\Im(z)$},
                xtick={-2,0,2,4,6,8,10}
            ]

            % Interior of shaded area imported as picture since it is beyond
            % TikZ fillbetween's capabilities
%             \addplot[opacity=.5] graphics [
%                 xmin=-3, xmax=12,
%                 ymin=-3, ymax=3,
%             ] {\plotdir/abs_jsun_geq1.png};
            % Never mind -- join_contours.py brought it within those capabilities!
            \addplot[mark=none, draw=none, fill=realextra!40, even odd rule] table {\plotdir/re12.dat} -- cycle;
%             % Visualizing the effects of join_contours.py
%             \addplot[red, mark=o, draw=none] table {\plotdir/re12.dat.log};

            % All interesting points, with debugging visualization commented out
            % zim0NM is the Mth real zero of the Nth factor of the polynomial
            %  determining the imaginary part
            %  (the locus of all zeroes is in \plotdir/im0N.dat).
            % xingN(u,l) is the Nth complex crossing point of the curves
            %  in the (upper,lower) half-plane
            \begin{scope}[every node/.style={draw,circle,inner sep=0pt}]

                \coordinate (zim011) at (0,0);
                \coordinate (zim012) at (3,0);
%                     \node[plotII] at (zim011) {\tiny1};
%                     \node[plotII] at (zim012) {\tiny2};

                \coordinate (zim021) at ({-2*sqrt(3) + 3},0);
                \coordinate (zim022) at (1,0);
                \coordinate (zim023) at ({2*sqrt(3) + 3},0);
%                     \node[plotI] at (zim021) {\tiny1};
%                     \node[plotI] at (zim022) {\tiny2};
%                     \node[plotI] at (zim023) {\tiny3};

                \coordinate (zim031) at ({-sqrt(3)*sqrt(2*sqrt(3) + 3) + sqrt(3) + 3},0);
                \coordinate (zim032) at ({4^(2/3) + 2*4^(1/3) + 3},0);
                \coordinate (zim033) at ({sqrt(3)*sqrt(2*sqrt(3) + 3) + sqrt(3) + 3},0);
%                     \node[plotIV] at (zim031) {\tiny1};
%                     \node[plotIV] at (zim032) {\tiny2};
%                     \node[plotIV] at (zim033) {\tiny3};

                \coordinate (zre121) at (1.740622, 0); % Root of an irreducible 6th-degree polynomial
                \coordinate (zre122) at (8.969980, 0); % this one too
%                     \node[plotVI] at (zre121) {\tiny1};
%                     \node[plotVI] at (zre122) {\tiny2};

                \coordinate (xing1u) at ({-1/2*4^(2/3)*(1) - 4^(1/3)*(1) + 3},{-1/2*4^(2/3)*(sqrt(3)) - 4^(1/3)*(-sqrt(3))});
                \coordinate (xing1l) at ({-1/2*4^(2/3)*(1) - 4^(1/3)*(1) + 3},{1/2*4^(2/3)*(sqrt(3)) + 4^(1/3)*(-sqrt(3))});
                \coordinate (xing2u) at (1.26795,+1.17996); % Analytic solutions don't find this one
                \coordinate (xing2l) at (1.26795,-1.17996);
%                     \node at (xing1u) {\tiny1};
%                     \node at (xing2u) {\tiny2};
            \end{scope}

%                 \addplot[mark=none, plotVII, samples=2, domain=-3:12] {0};
%                 \addplot[mark=none, plotII] table {\plotdir/im01.dat} ;
%                 \addplot[mark=none, plotI] table {\plotdir/im02.dat};
%                 \addplot[mark=none, plotIV] table {\plotdir/im03.dat};

            % Piecewise curves usling clipped scopes
            \begin{scope}
                \clip (zim032) |- (12,.25) -- (12,3) -| (zim012) |- (12,-3) -- (12,-.25) -| cycle;
                \addplot[mark=none, realneg, nocut] table {\plotdir/im01.dat};
                \addplot[mark=none, realpos, nocut] table {\plotdir/im02.dat};
                \addplot[mark=none, realneg, nocut] table {\plotdir/im03.dat};
            \end{scope}
            \begin{scope}
                \clip (zim032) |- (12,.25) -- (12,-.25) -| cycle;
                \addplot[mark=none, realpos] table {\plotdir/im03.dat};
            \end{scope}
            \begin{scope}
                \clip (xing1u) -- (-3,2) -- (-3,-2) -- (xing1l) -- cycle;
                \addplot[mark=none, realneg] table {\plotdir/im01.dat};
%                     \addplot[mark=none, realpos, branchcut] table {\plotdir/im02.dat};
                % Plots aren't compatible with zigzag lines, so approximate it with a bezier
                \draw[realpos, branchcut] (xing1u) .. controls (-.65,+.55) and (-.65,-.55) .. (xing1l);
            \end{scope}
            \begin{scope}
                \clip (xing1u) -- (1,+1) -| (zim022) |- (1,-1) -- (xing1l) -- cycle;
                \addplot[mark=none, realneg] table {\plotdir/im02.dat};
                \addplot[mark=none, realpos] table {\plotdir/im03.dat};
            \end{scope}
            \begin{scope}
                \clip (xing1u) -- (1,+1) -| (zim022) |- (1,-1)  -- (xing1l) -- (-3,-3) -| (zim012) |- (-3,+3) -- cycle;
                \addplot[mark=none, realpos] table {\plotdir/im01.dat};
                \addplot[mark=none, realpos, nocut] table {\plotdir/im02.dat};
                \addplot[mark=none, realneg, nocut] table {\plotdir/im03.dat};
            \end{scope}

            % Cuts along the real axis
            \draw[realpos, nocut] (-3,0) -- (zim021);
            \draw[realpos, branchcut] (zim021) -- (zim022);
            \draw[realneg, branchcut] (zim022) -- (zim012);
            \draw[realpos, branchcut] (zim012) -- (zim032);
            \draw[realneg, nocut] (zim032) -- (9,0);
            \draw[realpos, nocut] (9,0) -- (12,0);

            % Monodromy arrows
            \begin{scope}
                \draw[C1] (8.9,0.05) to[bend left=10] (8.9,0.3);
                \draw[C1] (2.6,1) to[bend left=10] (2.2,.5);
                \draw[C1] (0.45,1.3) to[bend left=10] (0.7,.7);
                \draw[C0] (0.9,0.4) to[bend left=20] (0.5,0.2);
                \draw[C1] (0.45,0.35) to[bend left=10] (0.15,0.2);
                \draw[C0] (0.1,0.2) -- (-0.2,0.2);

                \draw[C1] (8.9,-0.05) to[bend right=10] (8.9,-0.3);
                \draw[C1] (2.6,-1) to[bend right=10] (2.2,-.5);
                \draw[C1] (0.45,-1.3) to[bend right=10] (0.7,-.7);
                \draw[C0] (0.5,-0.2) to[bend left=20] (0.9,-0.4);
                \draw[C1] (0.45,-0.35) to[bend right=10] (0.15,-0.2);
                \draw[C0] (-0.2,-0.2) -- (0.1,-0.2);
            \end{scope}

            % Omega boundary
            \draw[gray, densely dotted] (0,0) circle[radius=3];

            \node at (5,2) {$\star$};
        \end{axis}

    \end{tikzpicture}
    \figrhoFcaption
    \label{fig:rhoF}
\end{figure}

%%% For monodromy.tex %%%

\begin{figure}[tp]
    \centering
    \newcommand{\drawaxes}{%
        \draw[->] (-1,0) -- (2.1,0);
        \draw[->] (0,-1) -- (0,+1);
        \draw[thin] (1,-.05) -- (1,+.05);
        \node at (1,.15) {\tiny $1$};}
    \newcommand{\drawcontours}[1][]{%
        \draw[thick,#1] (0,0) circle[radius=.5];
        \draw[thick,-stealth,#1] (0,0) ++ (135:.5) node[above left] {$C_0$} -- +(.01,.01);
        \draw[thick,-stealth,#1] (1,0) circle[radius=.5];
        \draw[thick,-stealth,#1] (1,0) ++ (45:.5) node[above right] {$C_1$} -- +(.01,-.01);}

    \tikzsetnextfilename{fig-monodromy-left}
    \begin{tikzpicture}[scale=1.2]
        \drawaxes
        \draw[thick, plotV, branchcut] (1,0) -- (2,0);
        \drawcontours
    \end{tikzpicture}
    \quad
    \tikzsetnextfilename{fig-monodromy-mid}
    \begin{tikzpicture}[scale=1.2]
        \drawaxes
        \draw[thick, plotV, branchcut] (-1,0) -- (0,0);
        \drawcontours
    \end{tikzpicture}
    \quad
    \tikzsetnextfilename{fig-monodromy-right}
    \begin{tikzpicture}[scale=1.2]
        \drawaxes
        \draw[thick, plotV, branchcut] (0,0) -- (1,0);
        \drawcontours[dashed]
        \draw[thick] (0,.5) arc[start angle=90, end angle=270, radius=.5] -- (1,-.5) arc[start angle=270, end angle=450, radius=.5] -- cycle;
    \end{tikzpicture}
    \figmonodromycaption
    \label{fig:monodromy}
\end{figure}

%%% For the worked example %%%

\begin{figure}

    % Label for the example coordinate
    \newcommand{\ex}{\text{ex}}
    % Coordinates for tstar = -2
    \newcommand{\retaustar}{0.}
    \newcommand{\imtaustar}{\suprad}
    \newcommand{\rebetastar}{1.73205}
    \newcommand{\imbetastar}{0.}
    % Coordinates for t = 4.2
    \newcommand{\retau}{0.263069100239244}
    \newcommand{\imtau}{0.164915378518587}
    \newcommand{\rebeta}{0.218217890235993}
    \newcommand{\imbeta}{0.}
    % Corner of the "too close to cut" zone
    \newcommand{\recorner}{1.}
    \newcommand{\imcorner}{0.4225}

    \hspace{-.5cm}
    \tikzsetnextfilename{fig-numeric-example}
    \begin{tikzpicture}[
        contour/.style={black, fill=#1, decorate, decoration={triangles,segment length=5pt}},
        integral annotation/.style={draw, rectangle, anchor=west},
        every text node part/.style={align=left}
        ]
        % tau plane
        \begin{axis}[
                xshift=-1cm,
                clip=false, axis on top=false,
                width=5cm, height=5cm, scale only axis, clip=false,
                xmin=0, xmax=.5,
                ymin=0, ymax=.5,
                xlabel={$\Re(\tau)$},
                ylabel={$\Im(\tau)$},
                xtick={0,.25,.5},
                xticklabels={$0$,$\tfrac14$, $\tfrac12$},
                ytick={0,\suprad/2,\suprad,.5},
                yticklabels={$0$,$\tfrac{\sqrt3}{12}$,$\tfrac{\sqrt3}{6}$,$\tfrac12$},
                yticklabel pos=left
            ]

            % The principal inverse image of the real line
            \draw[realline, spacelike] (0,.5) -- (0,0);
            \draw[realline, subthr] (0,0)
                arc[start angle=180, end angle=60, radius=\subrad];
            \draw[realline, twopi] (.5,\suprad)
                arc[start angle=90, end angle=150, radius=\suprad];
            \draw[realline, fourpi] (.5,\suprad) -- (.5,.5);

            % Coordinates for contours
            \coordinate (tau contour start) at (\retau,\imtau);
            \coordinate (tau contour end) at (\retaustar,\imtaustar);
            \coordinate (tau contour mid) at ($(tau contour start)!.5!(tau contour end)$);

            \coordinate (tau real start) at (tau contour end);
            \coordinate (tau real end) at (0,.53);
            \coordinate (tau real mid) at ($(tau real start)!.5!(tau real end)$);
        \end{axis}

        % tau contours
        \draw[contour=black]
            (tau contour start) -- (tau contour end)
                node[pos=.6, sloped, below] {$\mathcal C_\tau(\beta_\ex)$};
        \draw[contour=white]
               (tau real start) -- (tau real end);

        % tau points
        \draw
            (tau contour start)
                % circle[radius=.1pt]
                node [pin={[pin distance=1ex]270:{\footnotesize$
                    \begin{aligned}
                        \tau_\ex
                            &= 0.263+0.165i\\
                            &= \tfrac12 + \tfrac{\sqrt3}{6}e^{0.806i\pi}
                    \end{aligned}$}}] {}
            (tau contour end)
                % circle[radius=.1pt]
                node [pin=45:{$\tau_\star=\tfrac{i\sqrt3}{6}$}] {};

        % tau annotations
        \node[integral annotation] (IE) at (-2,7.5) {\footnotesize$
            \begin{alignedat}{2}
                \IE1(\beta_\star) &= (-6.020)_{<\chi} - (0.095)_{>\chi} &&= -6.115\\\IE2(\beta_\star) &= (\phantom{+}2.826)_{<\chi} + (0.001)_{>\chi} &&= +2.827
            \end{alignedat}
            $};
        \draw[-latex] (IE.west) .. controls ($(IE.west)+(-1,-1)$) and ($(tau real mid)+(-2,0)$) .. (tau real mid);
        \node[integral annotation] (Ctau) at (-1,6) {\footnotesize$
            \begin{aligned}
                \textstyle\int_{C_\tau(\beta_\ex)} h_1\ldots\d\tau &= 44.172+16.284i\\
                \textstyle\int_{C_\tau(\beta_\ex)} h_2\ldots\d\tau &= \phantom{0}8.462-67.465i
            \end{aligned}
            $};
        \draw[-latex] (Ctau.south) to[bend left=30] (tau contour mid);

        % xi plane
        \begin{axis}[
                xshift=6cm,
                clip=false, axis on top=false,
                width=5cm, height=5cm, scale only axis,
                xmin=0, xmax=2.5,
                ymin=0, ymax=1.5,
                xlabel={$\Re(\xi)$},
                ylabel={$\Im(\xi)$},
                xtick={0,1},
                xticklabels={$0$,$1$},
                ytick={0,1},
                yticklabels={$0$,$1$}
                ]
            ]

            % "Too close to cut" zone
            \fill[shaded] (0,0) -- (\recorner,0) -- (\recorner,\imcorner) -- (0,1) -- cycle;

            % The principal inverse image of the real line
            \draw[spacelike, realline] (+1,0) -- (+2.5,0);
            \draw[subthr, realline] (0,0) -- (0,1.5);
            \draw[twopi, realline] (0,0) -- ({+sqrt(3)/2},0);
            \draw[fourpi, realline] ({+sqrt(3)/2},0) -- (+1,0);

            % Coordinates for contours
            \coordinate (xi contour start) at (\rebeta,\imbeta);
            \coordinate (xi contour corner) at (\recorner,\imcorner);
            \coordinate (xi contour end) at (\rebetastar,\imbetastar);
            \coordinate (xi contour mid) at ($(xi contour corner)!.5!(xi contour end)$);

            \coordinate (xi real start) at (xi contour end);
            \coordinate (xi real end) at (2.8,0);
            \coordinate (xi real mid) at ($(xi real start)!.5!(xi real end)$);
        \end{axis}

        % xi contours
        \draw[contour=black]
               (xi contour start)
            -- (xi contour corner)
                node [midway, sloped, above] {$\mathcal C_\xi(\beta)$}
            -- (xi contour end);
        \draw[contour=white]
               (xi real start) -- (xi real end);

        % xi points
        \draw
           (xi contour start)
            % circle[radius=2pt]
            node[pin=-90:{$\beta_\ex=0.218$}] {}
           (xi contour end)
            % circle[radius=2pt]
            node[pin=-60:{$\beta_\star=\sqrt{3}$}] {};

        % xi annotations
        \node[integral annotation] (Cxi) at (6.5,6) {\footnotesize$
            \begin{aligned}
                \textstyle\int_{\mathcal C_\xi(\beta_\ex)}g_1\ldots\d\xi &= 70.061 - 21.617i\\
                \textstyle\int_{\mathcal C_\xi(\beta_\ex)}g_2\ldots\d\xi &= \phantom{0}2.283 - 45.262i
            \end{aligned}
            $};
        \draw[-latex] (Cxi.south) to[bend left=30] (xi contour mid);
        \node[integral annotation] (IJ) at (5.5,7.5) {\footnotesize$
            \begin{alignedat}{2}
                \IJ1(\beta_\star) &= (-0.897)_{<\chi} - (2\cdot10^{-4\phantom{0}})_{>\chi} &&= -0.897\\
                \IJ2(\beta_\star) &= (\phantom{+}0.008)_{<\chi} + (1\cdot10^{-17})_{>\chi} &&= 0.008
            \end{alignedat}
            $};
        \draw[-latex] (IJ.east) .. controls ($(IJ.east)+(+1,-1)$) and ($(xi real mid)+(2,2)$) .. (xi real mid);

        % beta annotations
        \node[integral annotation] (Irat) at (0,9) {\footnotesize$
            \begin{aligned}
                \Irat1(\beta_\ex) &= 21.034 - 2.937i\\
                \Irat2(\beta_\ex) &= \phantom{0}2.394 - 1.602i
            \end{aligned}$
            \qquad $\Idiv(\beta_\ex)  = 47.842 - 8.179i$\quad};

        % complete integrals
        \node[integral annotation] (E5) at (-1.5,-3) {\footnotesize$
            \begin{aligned}
                G_{15}(\beta_\ex) &= \phantom{0}{-0.952}\\
                G_{25}(\beta_\ex) &= \phantom{0}{-0.215 + \phantom{0}0.748i}\\
                B_{5}(\beta_\ex)  &= \phantom{-}12.853\\
                \I[\beta_\ex; G_{15}(\beta_\ex), G_{25}(\beta_\ex)] &= -10.294 + 18.144i\\
                \boldsymbol{\bar E_5(4,t_\ex)} &\boldsymbol{= 2.559 + 18.144i}
            \end{aligned}
        $};
        \node[integral annotation] (E6) at (5.5,-3) {\footnotesize$
            \begin{aligned}
                G_{16}(\beta_\ex) &= \phantom{-0}0\\
                G_{26}(\beta_\ex) &= \phantom{-0}0.546\\
                B_{6}(\beta_\ex)  &= -48.299 + 0.841i\\
                \I[\beta_\ex; G_{16}(\beta_\ex), G_{26}(\beta_\ex)] &= \phantom{-}96.069 - 8.974i\\
                \boldsymbol{\bar E_6(4,t_\ex)} &\boldsymbol{= 47.770 - 8.134i}
            \end{aligned}
        $};
    \end{tikzpicture}
    \fignumericexamplecaption
    \label{fig:numeric-example}
\end{figure}

%%% For the hybrid integral analysis %%%

\begin{figure}
    \hspace{-1.5cm}

    \def\hybridfile{\plotdir/hybrid_err_t-2.0}
    \input{\hybridfile_Hreg_reldiff_symlog.tex}

    \tikzset{
         SJ1/.style={plotI,   solid,              very thick},
         SJ2/.style={plotII,  densely dashed,     very thick},
        Hreg/.style={plotIII, densely dotted,     very thick},
        E1h2/.style={plotIV,  densely dashdotted, very thick}
    }

    \tikzsetnextfilename{fig-hybrid}
    \begin{tikzpicture}
        \begin{semilogxaxis}[
            width=.8\textwidth, height=\Eplotheight, scale only axis,
            xmin=3, xmax=1e8,
            xmajorgrids=true,
%             xlabel={$\chi$},
            xticklabels={},
            ymin=\symlogymin, ymax=\symlogymax,
            ytick={\symlogsparseytick},
            yticklabels={$-10^{-3}$,$-10^{-6}$,$-10^{-9}$,$-10^{-12}$,,$0$,,$-10^{-12}$,$+10^{-9}$,$+10^{-6}$,$+10^{-3}$},
            yticklabel style={font=\tiny},
            ymajorgrids=true,
            legend pos=outer north east
            ]

            \addplot[SJ1]  table[x=chi, y=result, col sep=tab] {\hybridfile_SJ1_reldiff_symlog.dat};
            \addplot[SJ2]  table[x=chi, y=result, col sep=tab] {\hybridfile_SJ2_reldiff_symlog.dat};
            \addplot[Hreg] table[x=chi, y=result, col sep=tab] {\hybridfile_Hreg_reldiff_symlog.dat};
            \addplot[E1h2] table[x=chi, y=result, col sep=tab] {\hybridfile_E1h2_reldiff_symlog.dat};

            \legend{$\IJ1$,$\IJ2$,$\IE1$,$\IE2$}
        \end{semilogxaxis}
        \begin{loglogaxis}[
            shift={(0,-(\Eplotheight + \Eplotsep)},
            width=.8\textwidth, height=\Eplotheight, scale only axis,
            xmin=1, xmax=1e8,
            xmajorgrids=true,
            xlabel={$\chi$},
            ymajorgrids=true,
            ]

            \addplot[SJ1]  table[x=chi, y=dresult, col sep=tab] {\hybridfile_SJ1.dat};
            \addplot[SJ2]  table[x=chi, y=dresult, col sep=tab] {\hybridfile_SJ2.dat};
            \addplot[Hreg] table[x=chi, y=dresult, col sep=tab] {\hybridfile_Hreg.dat};
            \addplot[E1h2] table[x=chi, y=dresult, col sep=tab] {\hybridfile_E1h2.dat};

%             \addplot[SJ1,  thin] table[x=chi, y=dhead, col sep=tab] {\hybridfile_SJ1.dat};
%             \addplot[SJ2,  thin] table[x=chi, y=dhead, col sep=tab] {\hybridfile_SJ2.dat};
%             \addplot[Hreg, thin] table[x=chi, y=dhead, col sep=tab] {\hybridfile_Hreg.dat};
%             \addplot[E1h2, thin] table[x=chi, y=dhead, col sep=tab] {\hybridfile_E1h2.dat};
        \end{loglogaxis}
        \useasboundingbox (0,-1.2\Eplotheight) + (0,-\Eplotsep) + (0,-.5cm); % Force larger bounding box
    \end{tikzpicture}

    \fighybridcaption
    \label{fig:hybrid}
\end{figure}

%%% For the result plots %%%

\begin{figure}
    \centering
    \tikzsetnextfilename{fig-Jbub}
    \begin{tikzpicture}
        \begin{axis}[
            width=.8\textwidth, height=1.5*\Eplotheight, scale only axis,
            xmin=\Eplotxmin, xmax=\Eplotxmax,
            ymin=-4.5, ymax=+4.5,
            xlabel={$t$},
            ylabel={$\frac{1}{n!}\Jbub{n}(t)$},
            xtick={0,4,16}, xmajorgrids=true,
            minor xtick={-9,-8,...,23},
            legend pos=north west, reverse legend
            ]

            \plotJbub{3}
            \plotJbub{2}
            \plotJbub{1}

            \legend{$n=3$,$n=2$,$n=1$}
        \end{axis}
    \end{tikzpicture}
    \figJbubcaption
    \label{fig:Jbub}
\end{figure}

\begin{figure}
    \hspace{-1cm}
    \tikzsetnextfilename{fig-masters-E1}
    \plotE{1}{E_1(2;t)}{-24}{+4}{16}
    \tikzsetnextfilename{fig-masters-E4}
    \plotE{4}{\bar E_4(4;t)}{-32}{+32}{4}

    \hspace{-1cm}
    \tikzsetnextfilename{fig-masters-E2}
    \plotE{2}{E_2(2;t)}{-4}{+24}{16}
    \tikzsetnextfilename{fig-masters-E5}
    \plotE{5}{\bar E_5(4;t)}{-32}{+32}{4}

    \hspace{-1cm}
    \tikzsetnextfilename{fig-masters-E3}
    \plotE{3}{E_3(2;t)}{-16}{+16}{16}
    \tikzsetnextfilename{fig-masters-E6}
    \plotE{6}{\bar E_6(4;t)}{-32}{+32}{4}

    \hspace{3cm}$t$\hspace{6.5cm}$t$ % Very hacky but gets the job done

    \figmasterscaption
    \label{fig:masters}
\end{figure}

\begin{figure}
    \hspace{-1cm}
    \tikzsetnextfilename{fig-reldiff-E1}
    \plotEsymlog{1}{E_1(2;t)}{16}
    \tikzsetnextfilename{fig-reldiff-E4}
    \plotEsymlog{4}{\bar E_4(4;t)}{4}

    \hspace{-1cm}
    \tikzsetnextfilename{fig-reldiff-E2}
    \plotEsymlog{2}{E_2(2;t)}{16}
    \tikzsetnextfilename{fig-reldiff-E5}
    \plotEsymlog{5}{\bar E_5(4;t)}{4}

    \hspace{-1cm}
    \tikzsetnextfilename{fig-reldiff-E3}
    \plotEsymlog{3}{E_3(2;t)}{16}
    \tikzsetnextfilename{fig-reldiff-E6}
    \plotEsymlog{6}{\bar E_6(4;t)}{4}

    \hspace{2.7cm}$t$\hspace{6.8cm}$t$ % Very hacky but gets the job done

    \figreldiffcaption
    \label{fig:reldiff}
\end{figure}

\begin{figure}
    \hspace{-1cm}
    \tikzsetnextfilename{fig-time}
    \begin{tikzpicture}
        \begin{semilogyaxis}[
            width=.8\textwidth, height=2*\Eplotheight, scale only axis,
            xmin=\Eplotxmin, xmax=\Eplotxmax, xlabel={$t$},
            xtick={0,4,16},
            minor xtick={-9,-8,...,23},
            ymin=2e-3, ymax=25, ylabel={Time [seconds/sample]},
            ymajorgrids=true, xmajorgrids=true,
            legend pos=outer north east, legend cell align=left
            ]

            \plotTime{1}
            \plotTime{2}
            \plotTime{3}
            \plotTime{4}
            \plotTime{5}
            \plotTime{6}

            \legend{$E_1(2;t)$,$E_2(2;t)$,$E_3(2;t)$,$\bar E_4(4;t)$,$\bar E_5(4;t)$,$\bar E_6(4;t)$}

        \end{semilogyaxis}
    \end{tikzpicture}
    \figtimecaption
    \label{fig:time}
\end{figure}

\begin{figure}
    \def\Eplotfile{\plotdir/PiE.dat}
    \hspace{-1cm}
    \tikzsetnextfilename{fig-PiE}
    \begin{tikzpicture}
        \begin{axis}[
%                     width=\Eplotwidth,
                width=.8\textwidth,
                height=2*\Eplotheight, scale only axis,
%                     ylabel={$\Re\bar\Pi_E(t)$},
%                     ylabel={$\bar\Pi_E(t)$},
                xmin=\Eplotxmin, xmax=\Eplotxmax,
%                     xlabel={$t$},
                ylabel={$\bar\Pi(t)$},
                xtick={0,4,16}, xmajorgrids=true, xticklabel pos=top,
                minor xtick={-9,-8,...,23},
                ymin=-245, ymax=245, ytick distance=50,
                legend pos=outer north east, reverse legend, legend cell align=left
                ]

            %%% Real and imag together
            \addplot[real, complex legend={thick, gray}, thick, gray]                  table[x=t, y=eRePiE0, col sep=tab] {\Eplotfile};
            \addplot[imag, forget plot, thick, gray]                                   table[x=t, y=eImPiE0, col sep=tab] {\Eplotfile};
            \addplot[real, complex legend={thick, E6}, thick, E6]                      table[x=t, y=eRePiE6, col sep=tab] {\Eplotfile};
            \addplot[imag, forget plot, thick, E6, restrict x to domain=4:\Eplotxmax]  table[x=t, y=eImPiE6, col sep=tab] {\Eplotfile};
            \addplot[real, complex legend={thick, E5}, thick, E5]                      table[x=t, y=eRePiE5, col sep=tab] {\Eplotfile};
            \addplot[imag, forget plot, thick, E5, restrict x to domain=4:\Eplotxmax]  table[x=t, y=eImPiE5, col sep=tab] {\Eplotfile};
            \addplot[real, complex legend={thick, E3}, thick, E3]                      table[x=t, y=eRePiE3, col sep=tab] {\Eplotfile};
            \addplot[imag, forget plot, thick, E3, restrict x to domain=16:\Eplotxmax] table[x=t, y=eImPiE3, col sep=tab] {\Eplotfile};
            \addplot[real, complex legend={thick, E2}, thick, E2]                      table[x=t, y=eRePiE2, col sep=tab] {\Eplotfile};
            \addplot[imag, forget plot, thick, E2, restrict x to domain=16:\Eplotxmax] table[x=t, y=eImPiE2, col sep=tab] {\Eplotfile};
            \addplot[real, complex legend={thick, E1}, thick, E1]                      table[x=t, y=eRePiE1, col sep=tab] {\Eplotfile};
            \addplot[imag, forget plot, thick, E1, restrict x to domain=16:\Eplotxmax] table[x=t, y=eImPiE1, col sep=tab] {\Eplotfile};
            \addplot[real, series, black,          restrict x to domain=\Eplotxmin:16] table[x=t, y=0RePiE, col sep=tab] {\Eplotfile};
            \addplot[real, complex legend={total}, total]                              table[x=t, y=eRePiE, col sep=tab] {\Eplotfile};
            \addplot[imag, forget plot, total,     restrict x to domain=4:\Eplotxmax]  table[x=t, y=eImPiE, col sep=tab] {\Eplotfile};

            \legend{rem.,$\propto\bar E_6$,$\propto\bar E_5$,$\propto E_3$,$\propto E_2$,$\propto E_1$,series,$\bar \Pi_E$}

        \end{axis}
        \begin{semilogyaxis}[
                shift={(0,-(\Eplotheight + \Eplotsep))},
                width=.8\textwidth,
                height=\Eplotheight, scale only axis,
%                     ylabel={$\delta\bar\Pi_E(t)$},
                xmin=\Eplotxmin, xmax=\Eplotxmax,
                xtick={0,4,16}, xmajorgrids=true,
                minor xtick={-9,-8,...,23},
                xlabel={$t$},
                ylabel={$\delta\bar\Pi(t)$},
                ymin=1e-16, ymax=1e-5, ytick={1e-15,1e-12,1e-9,1e-6}, ymajorgrids=true, ytick style={draw=none},
                ]

            \addplot[real, total]         table[x=t, y=eAbsRePiE, col sep=tab]  {\Eplotfile};
            \addplot[imag, total]         table[x=t, y=eAbsImPiE, col sep=tab]  {\Eplotfile};
            \addplot[real, thick, gray]   table[x=t, y=eAbsRePiE0, col sep=tab] {\Eplotfile};
            \addplot[imag, thick, gray]   table[x=t, y=eAbsImPiE0, col sep=tab] {\Eplotfile};
            \addplot[real, thick, E6]     table[x=t, y=eAbsRePiE6, col sep=tab] {\Eplotfile};
            \addplot[imag, thick, E6]     table[x=t, y=eAbsImPiE6, col sep=tab] {\Eplotfile};
            \addplot[real, thick, E5]     table[x=t, y=eAbsRePiE5, col sep=tab] {\Eplotfile};
            \addplot[imag, thick, E5 ]    table[x=t, y=eAbsImPiE5, col sep=tab] {\Eplotfile};
            \addplot[real, thick, E3]     table[x=t, y=eAbsRePiE3, col sep=tab] {\Eplotfile};
            \addplot[imag, thick, E3]     table[x=t, y=eAbsImPiE3, col sep=tab] {\Eplotfile};
            \addplot[real, thick, E2]     table[x=t, y=eAbsRePiE2, col sep=tab] {\Eplotfile};
            \addplot[imag, thick, E2]     table[x=t, y=eAbsImPiE2, col sep=tab] {\Eplotfile};
            \addplot[real, thick, E1]     table[x=t, y=eAbsRePiE1, col sep=tab] {\Eplotfile};
            \addplot[imag, thick, E1]     table[x=t, y=eAbsImPiE1, col sep=tab] {\Eplotfile};
            \addplot[real, series, black] table[x=t, y=0AbsRePiE, col sep=tab]  {\Eplotfile};

        \end{semilogyaxis}
        \useasboundingbox (0,-1.2\Eplotheight) + (0,-\Eplotsep) + (0,-.5cm); % Force larger bounding box
    \end{tikzpicture}

    \figPiEcaption
    \label{fig:PiE}
\end{figure}


\begin{figure}
    \def\Jplotfile{\plotdir/PiJ.dat}
    \def\Zplotfile{\plotdir/PiZ.dat}
    \hspace{-1cm}
    \tikzsetnextfilename{fig-PiJ}
    \begin{tikzpicture}
        \begin{axis}[
                width=.8\textwidth,
                height=2*\Eplotheight, scale only axis,
%                     ylabel={$\Re\bar\Pi_J(t)$},
%                     ylabel={$\bar\Pi_J(t)$},
                xmin=\Eplotxmin, xmax=\Eplotxmax,
%                     xlabel={$t$},
                xtick={0,4,16}, xmajorgrids=true,
                minor xtick={-9,-8,...,23},
                xlabel={$t$},
                ylabel={$\bar\Pi(t)$},
                ymin=-345, ymax=345, ytick distance=50,
                legend pos=outer north east, reverse legend, legend cell align=left
                ]

            \addplot[real, total, seriesdotted]                                        table[x=t, y=eRePiZ, col sep=tab] {\Zplotfile};
            \addplot[real, complex legend={thick, gray}, thick, gray]                  table[x=t, y=eRePiJ0, col sep=tab] {\Jplotfile};
            \addplot[imag, forget plot, thick, gray]                                   table[x=t, y=eImPiJ0, col sep=tab] {\Jplotfile};
            \addplot[real, complex legend={thick, E6}, thick, E6]                      table[x=t, y=eRePiJ111, col sep=tab] {\Jplotfile};
            \addplot[imag, forget plot, thick, E6, restrict x to domain=4:\Eplotxmax]  table[x=t, y=eImPiJ111, col sep=tab] {\Jplotfile};
            \addplot[real, complex legend={thick, E5}, thick, E5]                      table[x=t, y=eRePiJ21, col sep=tab] {\Jplotfile};
            \addplot[imag, forget plot, thick, E5, restrict x to domain=4:\Eplotxmax]  table[x=t, y=eImPiJ21, col sep=tab] {\Jplotfile};
            \addplot[real, complex legend={thick, E4}, thick, E4]                      table[x=t, y=eRePiJ11, col sep=tab] {\Jplotfile};
            \addplot[imag, forget plot, thick, E4, restrict x to domain=4:\Eplotxmax]  table[x=t, y=eImPiJ11, col sep=tab] {\Jplotfile};
            \addplot[real, complex legend={thick, E3}, thick, E3]                      table[x=t, y=eRePiJ3, col sep=tab] {\Jplotfile};
            \addplot[imag, forget plot, thick, E3, restrict x to domain=4:\Eplotxmax]  table[x=t, y=eImPiJ3, col sep=tab] {\Jplotfile};
            \addplot[real, complex legend={thick, E2}, thick, E2]                      table[x=t, y=eRePiJ2, col sep=tab] {\Jplotfile};
            \addplot[imag, forget plot, thick, E2, restrict x to domain=4:\Eplotxmax]  table[x=t, y=eImPiJ2, col sep=tab] {\Jplotfile};
            \addplot[real, complex legend={thick, E1}, thick, E1]                      table[x=t, y=eRePiJ1, col sep=tab] {\Jplotfile};
            \addplot[imag, forget plot, thick, E1, restrict x to domain=4:\Eplotxmax]  table[x=t, y=eImPiJ1, col sep=tab] {\Jplotfile};
            \addplot[real, complex legend={total}, total]                              table[x=t, y=eRePiJ, col sep=tab] {\Jplotfile};
            \addplot[imag, forget plot, total,     restrict x to domain=4:\Eplotxmax]  table[x=t, y=eImPiJ, col sep=tab] {\Jplotfile};

            \legend{$\bar\Pi_\zeta$,rem.,$\propto[\Jbub1]^3$,$\propto\Jbub1\Jbub2$,$\propto[\Jbub1]^2$,$\propto\Jbub3$,$\propto\Jbub2$,$\propto\Jbub1$,$\bar\Pi_J$}

        \end{axis}
    \end{tikzpicture}
    \figPiJcaption
    \label{fig:PiJ}
\end{figure}


% \begin{figure}
%     \centering
%     \newcommand{\plotwidth}{.34\columnwidth}
%     \hspace{-.5cm}
%     \tikzsetnextfilename{fig-tau-err-arg}
%     \begin{tikzpicture}
%         \begin{axis}[
%             width=\plotwidth, height=\plotwidth, scale only axis, axis on top,
%             xmin=0, xmax=.5,
%             ymin=0, ymax=.5,
%             xlabel={$\Re(\tau)$},
%             ylabel={$\Im(\tau)$},
%             xtick={-.5,-.25,0,+.25,+.5},
%             xticklabels={$-\tfrac12$, $-\tfrac14$,$0$,$\tfrac14$,$\tfrac12$},
%             ytick={0,\suprad/2,\suprad,.5},
%             yticklabels={$0$,$\tfrac{\sqrt3}{12}$,$\tfrac{\sqrt3}{6}$,$\tfrac12$}
%             ]
%             \addplot graphics [
%                     xmin=0, xmax=.5,
%                     ymin=.01, ymax=.5,
%                 ] {\plotdir/tau_t_tau.png};
%             % Real line to check figure alignment and hide ugly edge
%             \draw[thick, gray, cap=round] (0,.5) -- (0,0)
%                 arc[start angle=180, end angle=60, radius=\subrad]
%                 (.5,.5) -- (.5,\suprad)
%                 arc[start angle=90, end angle=150, radius=\suprad];
%             % Arrow mentioned in the legend
%             \draw[gray, -latex] (.112,.115) -- +(-1.5ex,1.5ex);
%         \end{axis}
%     \end{tikzpicture}
%     \tikzsetnextfilename{fig-tau-err-abs}
%     \begin{tikzpicture}
%         \begin{axis}[
%             width=\plotwidth, height=\plotwidth, scale only axis, axis on top,
%             xmin=0, xmax=.5,
%             ymin=0, ymax=.5,
%             xlabel={$\Re(\tau)$},
% %             ylabel={$\Im(\tau)$},
%             xtick={-.5,-.25,0,+.25,+.5},
%             xticklabels={$-\tfrac12$, $-\tfrac14$,$0$,$\tfrac14$,$\tfrac12$},
%             ytick={0,\suprad/2,\suprad,.5},
%             yticklabels={},
%             ]
%             \addplot graphics [
%                     xmin=0, xmax=.5,
%                     ymin=.018, ymax=.5,
%                 ] {\plotdir/tau_t_tau_abs.png};
%             % Real line to check figure alignment and hide ugly edge
%             \draw[thick, gray, cap=round] (0,.5) -- (0,0)
%                 arc[start angle=180, end angle=60, radius=\subrad]
%                 (.5,.5) -- (.5,\suprad)
%                 arc[start angle=90, end angle=150, radius=\suprad];
%             % Arrow mentioned in the legend
%             \draw[gray, -latex] (.112,.115) -- +(-1.5ex,1.5ex);
%         \end{axis}
%     \end{tikzpicture}
%     \tikzsetnextfilename{fig-tau-err-cb}
%     \begin{tikzpicture}
%         \begin{semilogyaxis}[
%             width=.5cm, height=\plotwidth, scale only axis, axis on top,
%             xmin=0, xmax=1,
%             ymin=1e-23, ymax=1e-17,
%             xlabel={\vphantom{$\Re(\tau)$}},
%             ylabel={$|\delta\tau|$},
%             yticklabel pos=right,
%             xtick={.5}, xticklabels={\vphantom{$\tfrac12$}}, xtick style={draw=none}]
%
%             \addplot graphics [
%                     xmin=0, xmax=1,
%                     ymin=1e-23, ymax=1e-17,
%                 ] {\plotdir/tau_t_tau_cb.png};
%         \end{semilogyaxis}
%     \end{tikzpicture}
%     \figtauerrcaption
%     \label{fig:tau-err}
% \end{figure}

\end{document}
